\input{preamble}

% OK, start here.
%
\begin{document}

\title{Foncteurs et morphismes}


\maketitle

\phantomsection
\label{section-phantom}

\tableofcontents

\section{Introduction}
\label{section-introduction}

\noindent
Soient $X$ et $Y$ des schémas. Ce chapitre porte sur les rapports
entre les foncteurs $\QCoh(\mathcal{O}_Y) \to \QCoh(\mathcal{O}_X)$ et les
morphismes de schémas $X \to Y$. Plus généralement, nous étudions les
rapports entre $\QCoh(\mathcal{O}_X)$ et $X$ ou, si $X$ est noethérien,
entre $\textit{Coh}(\mathcal{O}_X)$ et $X$.
Ces rapports ont été étudiés dans \cite{Gabriel}.






\section{Foncteurs sur les catégories de modules}
\label{section-preliminary}

\noindent
Pour un anneau $A$, notons $\text{Mod}^{fp}_A$ la catégorie des
$A$-modules de présentation finie.

\begin{lemma}
\label{lemma-functor-on-fp-modules}
Soit $A$ un anneau. Soit $\mathcal{B}$ une catégorie admettant des limites
inductives filtrantes. Soit $F : \text{Mod}^{fp}_A \to \mathcal{B}$ un foncteur. Alors $F$
se prolonge de manière unique en un foncteur $F' : \text{Mod}_A \to \mathcal{B}$
qui commute aux limites inductives filtrantes.
\end{lemma}

\begin{proof}
Cela résulte de
Catégories, lemme \ref{categories-lemma-extend-functor-by-colim}.
Pour voir que ce lemme s'applique, observons que les
$A$-modules de présentation finie sont des
objets compacts au sens catégorique de $\text{Mod}_A$, d'après
Algèbre, lemme \ref{algebra-lemma-characterize-finitely-presented-module-hom}.
De plus, tout $A$-module est limite inductive filtrante
de $A$-modules de présentation finie, d'après
Algèbre, lemme \ref{algebra-lemma-module-colimit-fp}.
\end{proof}

\noindent
Si une catégorie $\mathcal{B}$ est additive et admet des limites inductives filtrantes,
alors $\mathcal{B}$ admet des sommes directes quelconques : toute somme directe s'écrit
comme limite inductive filtrante de sommes directes finies.

\begin{lemma}
\label{lemma-functor-on-fp-modules-additive}
Soient $A$, $\mathcal{B}$, $F$ comme dans le lemme \ref{lemma-functor-on-fp-modules}.
Supposons que $\mathcal{B}$ et $F$ soient additifs. Alors
$F'$ est additif et commute aux sommes directes quelconques.
\end{lemma}

\begin{proof}
Pour montrer que $F'$ est additif, il suffit de montrer
que $F'(M) \oplus F'(M') \to F'(M \oplus M')$ est un isomorphisme pour
tous $A$-modules $M$, $M'$, voir
Homologie, lemme \ref{homology-lemma-additive-functor}.
Écrivons $M = \colim_i M_i$ et $M' = \colim_j M'_j$ comme limites inductives filtrantes
de $A$-modules de présentation finie $M_i$. Alors
$F'(M) = \colim_i F(M_i)$, $F'(M') = \colim_j F(M'_j)$, et
\begin{align*}
F'(M \oplus M')
& =
F'(\colim_{i, j} M_i \oplus M'_j) \\
& =
\colim_{i, j} F(M_i \oplus M'_j) \\
& =
\colim_{i, j} F(M_i) \oplus F(M'_j) \\
& =
F'(M) \oplus F'(M')
\end{align*}
comme voulu. Pour montrer que $F'$ commute aux sommes directes, supposons
que $M = \bigoplus_{i \in I} M_i$. Alors
$M = \colim_{I' \subset I\text{ fini}} \bigoplus_{i \in I'} M_i$
est une limite inductive filtrante. On obtient
\begin{align*}
F'(M)
& =
\colim_{I' \subset I\text{ fini}}
F'(\bigoplus\nolimits_{i \in I'} M_i) \\
& =
\colim_{I' \subset I\text{ fini}}
\bigoplus\nolimits_{i \in I'} F'(M_i) \\
& =
\bigoplus\nolimits_{i \in I} F'(M_i)
\end{align*}
La deuxième égalité résulte de l'additivité de $F'$ déjà démontrée.
\end{proof}

\noindent
Si une catégorie $\mathcal{B}$ est additive, admet des limites inductives filtrantes et
des conoyaux, alors $\mathcal{B}$ admet des limites inductives quelconques, voir
la discussion ci-dessus et Catégories, lemme
\ref{categories-lemma-colimits-coproducts-coequalizers}.

\begin{lemma}
\label{lemma-functor-on-fp-modules-right-exact}
Soient $A$, $\mathcal{B}$, $F$ comme dans le lemme \ref{lemma-functor-on-fp-modules}.
Supposons que $\mathcal{B}$ soit additive, admette des conoyaux, et que $F$ soit exact à droite. Alors
$F'$ est additif, exact à droite et commute aux sommes directes quelconques.
\end{lemma}

\begin{proof}
Puisque $F$ est exact à droite, $F$ commute aux sommes de deux objets, qui sont
représentées par les sommes directes. Ainsi $F$ est additif, d'après
Homologie, lemme \ref{homology-lemma-additive-functor}.
Donc $F'$ est additif et commute aux sommes directes, d'après le
lemme \ref{lemma-functor-on-fp-modules-additive}.
Nous invitons le lecteur à démontrer lui-même que $F'$ est exact à droite
avant de lire la démonstration ci-dessous.

\medskip\noindent
Pour montrer que $F'$ est exact à droite, il suffit de montrer que $F'$ commute
aux conoyaux de doubles flèches, voir
Catégories, lemme \ref{categories-lemma-characterize-right-exact}.
Or, si $a, b : K \to L$ sont des homomorphismes de $A$-modules, le
conoyau de la double flèche formée de $a$ et $b$ est le conoyau de $a - b : K \to L$.
Soit donc $K \to L \to M \to 0$ une suite exacte
de $A$-modules. Nous devons montrer que, dans
$$
F'(K) \to F'(L) \to F'(M) \to 0
$$
la deuxième flèche est un conoyau de la première dans $\mathcal{B}$
(si $\mathcal{B}$ était abélienne, nous dirions que la suite ci-dessus
est exacte).
Écrivons $M = \colim_{i \in I} M_i$ comme limite inductive filtrante de
$A$-modules de présentation finie, voir
Algèbre, lemme \ref{algebra-lemma-module-colimit-fp}.
Posons $L_i = L \times_M M_i$. 
On obtient un système de suites exactes $K \to L_i \to M_i \to 0$ indexé par $I$.
Puisque les limites inductives commutent aux limites inductives, d'après
Catégories, lemme \ref{categories-lemma-colimits-commute},
et puisque les conoyaux sont des cas particuliers des conoyaux de doubles flèches,
il suffit de montrer que $F'(L_i) \to F(M_i)$ est un conoyau de
$F'(K) \to F'(L_i)$ dans $\mathcal{B}$ pour tout $i \in I$. Autrement dit, nous pouvons
supposer $M$ de présentation finie. Écrivons $L = \colim_{i \in I} L_i$
comme limite inductive filtrante de $A$-modules de présentation finie,
de telle sorte que chaque $L_i$ s'envoie surjectivement sur $M$.
Posons $K_i = K \times_L L_i$. On obtient un système de suites exactes courtes
$K_i \to L_i \to M \to 0$ indexé par $I$.
En répétant l'argument précédent, on se ramène à montrer que
$F(L_i) \to F(M_i)$ est un conoyau de
$F'(K) \to F(L_i)$ dans $\mathcal{B}$ pour tout $i \in I$.
Autrement dit, nous pouvons supposer que
$L$ et $M$ sont des $A$-modules de présentation finie.
Dans ce cas, le module $\Ker(L \to M)$ est de type fini
(Algèbre, lemme \ref{algebra-lemma-extension}).
On peut donc écrire $K = \colim_{i \in I} K_i$ comme limite inductive filtrante
de $A$-modules de présentation finie s'envoyant tous surjectivement sur $\Ker(L \to M)$.
On obtient un système de suites exactes courtes
$K_i \to L \to M \to 0$ indexé par $I$.
En répétant l'argument précédent, on se ramène à montrer que
$F(L) \to F(M)$ est un conoyau de
$F(K_i) \to F(L)$ dans $\mathcal{B}$ pour tout $i \in I$.
Autrement dit, nous pouvons supposer que $K$, $L$ et $M$
sont des $A$-modules de présentation finie. Ce dernier cas résulte
de l'hypothèse que $F$ est exact à droite.
\end{proof}

\noindent
Si une catégorie $\mathcal{B}$ est additive et admet des noyaux,
alors $\mathcal{B}$ admet des limites projectives finies. En effet, les produits finis
sont des sommes directes, qui existent, et le noyau de la double flèche $a, b : L \to M$
est le noyau de $a - b : K \to L$, qui existe. Toutes les limites projectives
finies existent donc, d'après Catégories, lemme \ref{categories-lemma-finite-limits-exist}.

\begin{lemma}
\label{lemma-functor-on-fp-modules-left-exact}
Soient $A$, $\mathcal{B}$, $F$ comme dans le lemme \ref{lemma-functor-on-fp-modules}.
Supposons que $A$ soit un anneau cohérent (Algèbre, définition
\ref{algebra-definition-coherent}), que $\mathcal{B}$ soit additive, admette des noyaux,
que les limites inductives filtrantes commutent à la formation des noyaux et que $F$ soit exact à gauche. Alors
$F'$ est additif, exact à gauche et commute aux sommes directes quelconques.
\end{lemma}

\begin{proof}
Puisque $A$ est cohérent, la catégorie $\text{Mod}^{fp}_A$ est abélienne,
avec les mêmes noyaux et conoyaux que dans $\text{Mod}_A$, voir
Algèbre, lemmes \ref{algebra-lemma-coherent-ring} et
\ref{algebra-lemma-coherent}. Toutes les limites projectives finies existent donc dans
$\text{Mod}^{fp}_A$, et
Catégories, définition \ref{categories-definition-exact}, s'applique.
Puisque $F$ est exact à gauche, $F$ commute aux produits de deux objets, qui sont
représentés par les sommes directes. Ainsi $F$ est additif, d'après
Homologie, lemme \ref{homology-lemma-additive-functor}.
Donc $F'$ est additif et commute aux sommes directes, d'après le
lemme \ref{lemma-functor-on-fp-modules-additive}.
Nous invitons le lecteur à démontrer lui-même que $F'$ est exact à gauche
avant de lire la démonstration ci-dessous.

\medskip\noindent
Pour montrer que $F'$ est exact à gauche, il suffit de montrer que $F'$ commute
aux noyaux de doubles flèches, voir
Catégories, lemme \ref{categories-lemma-characterize-left-exact}.
Or, si $a, b : L \to M$ sont des homomorphismes de $A$-modules, le
noyau de la double flèche formée de $a$ et $b$ est le noyau de $a - b : L \to M$.
Soit donc $0 \to K \to L \to M$ une suite exacte
de $A$-modules. Nous devons montrer que, dans
$$
0 \to F'(K) \to F'(L) \to F'(M)
$$
la flèche $F'(K) \to F'(L)$ est un noyau de $F'(L) \to F'(M)$ dans $\mathcal{B}$
(si $\mathcal{B}$ était abélienne, nous dirions que la suite ci-dessus
est exacte).
Écrivons $M = \colim_{i \in I} M_i$ comme limite inductive filtrante de
$A$-modules de présentation finie, voir
Algèbre, lemme \ref{algebra-lemma-module-colimit-fp}.
Posons $L_i = L \times_M M_i$. 
On obtient un système de suites exactes $0 \to K \to L_i \to M_i$
indexé par $I$. Puisque, par hypothèse, les limites inductives filtrantes commutent à la formation des noyaux
dans $\mathcal{B}$, 
il suffit de montrer que $F'(K) \to F'(L_i)$ est un noyau de
$F'(L_i) \to F(M_i)$ dans $\mathcal{B}$ pour tout $i \in I$. Autrement dit, nous pouvons
supposer $M$ de présentation finie. Écrivons $L = \colim_{i \in I} L_i$
comme limite inductive filtrante de $A$-modules de présentation finie.
Posons $K_i = K \times_L L_i$. On obtient un système de suites exactes courtes
$0 \to K_i \to L_i \to M$ indexé par $I$.
En répétant l'argument précédent, on se ramène à montrer que
$F'(K_i) \to F(L_i)$ est un noyau de
$F(L_i) \to F(M)$ dans $\mathcal{B}$ pour tout $i \in I$.
Autrement dit, nous pouvons supposer que
$L$ et $M$ sont des $A$-modules de présentation finie.
Puisque $A$ est cohérent, le $A$-module $K = \Ker(L \to M)$ est
de présentation finie, car la catégorie des
$A$-modules de présentation finie est abélienne (voir les références ci-dessus).
Autrement dit, les trois modules $K$, $L$ et $M$
sont des $A$-modules de présentation finie. Ce dernier cas résulte
de l'hypothèse que $F$ est exact à gauche.
\end{proof}

\noindent
Si une catégorie $\mathcal{B}$ est additive et admet des conoyaux,
alors $\mathcal{B}$ admet des limites inductives finies. En effet, les sommes finies
sont des sommes directes, qui existent, et le conoyau de la double flèche $a, b : K \to L$
est le conoyau de $a - b : K \to L$, qui existe. Toutes les limites inductives
finies existent donc, d'après Catégories, lemme \ref{categories-lemma-colimits-exist}.

\begin{lemma}
\label{lemma-functor-on-modules-fp}
Soit $A$ un anneau. Soit $\mathcal{B}$ une catégorie additive
admettant des conoyaux. Il existe une équivalence de catégories entre
\begin{enumerate}
\item la catégorie des foncteurs $F : \text{Mod}^{fp}_A \to \mathcal{B}$
exacts à droite, et
\item la catégorie des couples $(K, \kappa)$ où $K \in \Ob(\mathcal{B})$
et où $\kappa : A \to \text{End}_\mathcal{B}(K)$ est un homomorphisme d'anneaux,
\end{enumerate}
donnée par la règle qui à $F$ associe $F(A)$ muni de son action naturelle de $A$.
\end{lemma}

\begin{proof}
Soit $(K, \kappa)$ comme dans (2). Nous allons construire un foncteur
$F : \text{Mod}^{fp}_A \to \mathcal{B}$ tel que $F(A) = K$,
muni de l'action donnée de $A$ par $\kappa$. Pour un
entier $n \geq 0$, posons
$$
F(A^{\oplus n}) = K^{\oplus n}
$$
Étant donnée une application $A$-linéaire $\varphi : A^{\oplus m} \to A^{\oplus n}$
de matrice $(a_{ij}) \in \text{Mat}(n \times m, A)$, définissons
$$
F(\varphi) :
F(A^{\oplus m}) = K^{\oplus m}
\longrightarrow
K^{\oplus n} = F(A^{\oplus n})
$$
comme l'application de matrice $(\kappa(a_{ij}))$. Cela définit un foncteur
additif $F$ de la sous-catégorie pleine de
$\text{Mod}^{fp}_A$ d'objets $0$, $A$, $A^{\oplus 2}$, $\ldots$
dans $\mathcal{B}$ ; nous omettons la vérification.

\medskip\noindent
Pour chaque objet $M$ de $\text{Mod}^{fp}_A$, choisissons une présentation
$$
A^{\oplus m_M} \xrightarrow{\varphi_M} A^{\oplus n_M} \to M \to 0
$$
de $M$ comme $A$-module. Prenons la présentation triviale
$0 \to A^{\oplus n} \xrightarrow{1} A^{\oplus n} \to 0$ si $M = A^{\oplus n}$
(ce n'est pas nécessaire, mais simplifie l'exposé).
Pour chaque morphisme $f : M \to N$ de
$\text{Mod}^{fp}_A$, nous pouvons choisir un diagramme commutatif
\begin{equation}
\label{equation-map}
\vcenter{
\xymatrix{
A^{\oplus m_M} \ar[r]_{\varphi_M} \ar[d]_{\psi_f} &
A^{\oplus n_M} \ar[r] \ar[d]_{\chi_f} &
M \ar[r] \ar[d]_f & 0 \\
A^{\oplus m_N} \ar[r]^{\varphi_N} &
A^{\oplus n_N} \ar[r] &
N \ar[r] & 0
}
}
\end{equation}
Ces choix étant faits, nous pouvons définir ce qui suit : pour un objet
$M$ de $\text{Mod}^{fp}_A$, posons
$$
F(M) = \Coker(F(\varphi_M) : F(A^{\oplus m_M}) \to F(A^{\oplus n_M}))
$$
et, pour un morphisme $f : M \to N$ de $\text{Mod}^{fp}_A$, posons
$$
F(f) = \text{le morphisme }F(M) \to F(N)\text{ induit par }
F(\psi_f)\text{ et }F(\chi_f)\text{ sur les conoyaux}
$$
Notons que cette règle prolonge le foncteur $F$ donné sur
la sous-catégorie pleine constituée des modules libres $A^{\oplus n}$.
Il reste à montrer que $F$ est un foncteur, que $F$ est additif
et que $F$ est exact à droite.

\medskip\noindent
Soit $f : M \to N$ un morphisme de $\text{Mod}^{fp}_A$. Montrons que le morphisme
$F(f)$ définie ci-dessus est indépendante du choix de $\psi_f$ et de $\chi_f$
dans (\ref{equation-map}). Supposons en effet que
$$
\xymatrix{
A^{\oplus m_M} \ar[r]_{\varphi_M} \ar[d]_\psi &
A^{\oplus n_M} \ar[r] \ar[d]_\chi &
M \ar[r] \ar[d]_f & 0 \\
A^{\oplus m_N} \ar[r]^{\varphi_N} &
A^{\oplus n_N} \ar[r] &
N \ar[r] & 0
}
$$
soit également commutatif. Notons $F(f)' : F(M) \to F(N)$ le morphisme
induit par $F(\psi)$ et $F(\chi)$. À la vue des diagrammes
commutatifs, un argument élémentaire d'algèbre commutative fournit un homomorphisme
$\omega : A^{\oplus n_M} \to A^{\oplus m_N}$ telle que
$\chi = \chi_f + \varphi_N \circ \omega$. En appliquant $F$,
on trouve $F(\chi) = F(\chi_f) + F(\varphi_N) \circ F(\omega)$.
Puisque $F(N)$ est le conoyau de $F(\varphi_N)$, on voit
que le morphisme $F(A^{\oplus n_M}) \to F(M)$ égalise $F(f)$ et $F(f)'$.
Un conoyau étant un épimorphisme, on en conclut que $F(f) = F(f)'$.

\medskip\noindent
Montrons que $F$ est un foncteur. Observons d'abord que
$F(\text{id}_M) = \text{id}_{F(M)}$, car nous pouvons choisir
les identités pour $\psi_f$ et $\chi_f$ dans le diagramme ci-dessus
lorsque $f = \text{id}_M$. Supposons ensuite donnés
$f : M \to N$ et $g : L \to M$. Alors
$\psi = \psi_f \circ \psi_g$ et $\chi = \chi_f \circ \chi_g$
conviennent dans (\ref{equation-map}) pour $f \circ g$.
Ils induisent donc le morphisme voulu, ce qui exprime précisément
que $F(f) \circ F(g) = F(f \circ g)$.

\medskip\noindent
Montrons que $F$ est additif. Supposons donnés
$f, g : M \to N$. Alors $\psi = \psi_f + \psi_g$ et
$\chi = \chi_f + \chi_g$ conviennent dans (\ref{equation-map}) pour $f + g$.
Ils induisent donc le morphisme voulu, ce qui exprime précisément
que $F(f) + F(g) = F(f + g)$.

\medskip\noindent
Montrons enfin que $F$ est exact à droite. Il suffit de montrer que $F$
commute aux conoyaux de doubles flèches, voir
Catégories, lemme \ref{categories-lemma-characterize-right-exact}.
Pour cela, il suffit de montrer que $F$ commute aux conoyaux.
Soit $K \to L \to M \to 0$ une suite exacte de $A$-modules
avec $K$, $L$, $M$ de présentation finie. Puisque $F$ est un foncteur
additif, on obtient bien un complexe
$$
F(K) \to F(L) \to F(M) \to 0
$$
et il faut montrer que la deuxième flèche est le conoyau de la première
dans $\mathcal{B}$. En tout cas, on obtient un morphisme
$\Coker(F(K) \to F(L)) \to F(M)$.
Un argument élémentaire d'algèbre commutative fournit un diagramme commutatif
$$
\xymatrix{
A^{\oplus m_M} \ar[r]_{\varphi_M} \ar[d]_\psi &
A^{\oplus n_M} \ar[r] \ar[d]_\chi &
M \ar[r] \ar[d]_1 & 0 \\
K \ar[r] &
L \ar[r] &
M \ar[r] & 0
}
$$
En appliquant $F$ à ce diagramme et en utilisant la construction de $F(M)$ comme
conoyau de $F(\varphi_M)$, on trouve un morphisme
$F(M) \to \Coker(F(K) \to F(L))$ inverse à droite
du morphisme $\Coker(F(K) \to F(L)) \to F(M)$. Il en résulte d'abord
que $F(L) \to F(M)$ est toujours un épimorphisme. Ensuite, ce qui précède donne
une décomposition
$$
\Coker(F(K) \to F(L)) = F(M) \oplus E
$$
compatible à la fois avec
$F(M) \to \Coker(F(K) \to F(L))$ et $\Coker(F(K) \to F(L)) \to F(M)$.
Mais alors l'épimorphisme $p : F(L) \to E$ devient nul aussi bien
après composition avec $F(K) \to F(L)$ qu'après composition
avec $F(A^{n_M}) \to F(L)$. Or, puisque $K \oplus A^{n_M} \to L$
est surjectif (nous omettons l'argument algébrique), on en conclut que
$F(K \oplus A^{n_M}) \to F(L)$ est un épimorphisme (d'après ce qui précède),
d'où $E = 0$. Cela achève la démonstration.
\end{proof}

\begin{lemma}
\label{lemma-functor-on-modules}
Soit $A$ un anneau. Soit $\mathcal{B}$ une catégorie additive
admettant des sommes directes quelconques et des conoyaux. Il existe une équivalence
de catégories entre
\begin{enumerate}
\item la catégorie des foncteurs $F : \text{Mod}_A \to \mathcal{B}$
exacts à droite et commutant aux sommes directes quelconques, et
\item la catégorie des couples $(K, \kappa)$ où $K \in \Ob(\mathcal{B})$
et où $\kappa : A \to \text{End}_\mathcal{B}(K)$ est un homomorphisme d'anneaux,
\end{enumerate}
donnée par la règle qui à $F$ associe $F(A)$ muni de son action naturelle de $A$.
\end{lemma}

\begin{proof}
Combinons les lemmes \ref{lemma-functor-on-modules-fp} et
\ref{lemma-functor-on-fp-modules-right-exact}.
\end{proof}






\section{Foncteurs entre catégories de modules}
\label{section-functors}

\noindent
Le lemme suivant est l'archétype des résultats de ce chapitre.

\begin{lemma}
\label{lemma-functor}
Soient $A$ et $B$ des anneaux. Soit $F : \text{Mod}_A \to \text{Mod}_B$
un foncteur. Les propriétés suivantes sont équivalentes :
\begin{enumerate}
\item $F$ est isomorphe au foncteur $M \mapsto M \otimes_A K$
pour un certain $A \otimes_\mathbf{Z} B$-module $K$ ;
\item $F$ est exact à droite et commute à toutes les sommes directes ;
\item $F$ commute à toutes les limites inductives ;
\item $F$ admet un adjoint à droite $G$.
\end{enumerate}
\end{lemma}

\begin{proof}
Si (1) est satisfaite, alors (4) l'est aussi, car un adjoint à droite de $M \mapsto M \otimes_A K$
est $N \mapsto \Hom_B(K, N)$, voir
Algèbre différentielle graduée, lemme \ref{dga-lemma-tensor-hom-adjunction}.
Si (4) est satisfaite, alors (3) l'est d'après Catégories, lemme \ref{categories-lemma-adjoint-exact}.
L'implication (3) $\Rightarrow$ (2) résulte immédiatement des définitions.

\medskip\noindent
Supposons (2). Nous allons démontrer (1). D'après la discussion de
Homologie, section \ref{homology-section-functors},
le foncteur $F$ est additif. Ainsi, $F$ induit
un homomorphisme d'anneaux $A \to \text{End}_B(F(M))$, $a \mapsto F(a \cdot \text{id}_M)$
pour tout $A$-module $M$. Nous en concluons que $F(M)$ est un
$A \otimes_\mathbf{Z} B$-module, fonctoriellement en $M$.
Posons $K = F(A)$. Définissons
$$
M \otimes_A K = M \otimes_A F(A) \longrightarrow F(M),
\quad m \otimes k \longmapsto F(\varphi_m)(k)
$$
Ici, $\varphi_m : A \to M$ envoie $a \to am$. La règle
$(m, k) \mapsto F(\varphi_m)(k)$ est $A$-bilinéaire (et $B$-linéaire
à droite), comme requis pour obtenir l'homomorphisme
$A \otimes_\mathbf{Z} B$-linéaire affichée.
Cette construction est fonctorielle en $M$ ; elle définit donc une transformation
de foncteurs $- \otimes_A K \to F(-)$ qui est un isomorphisme lorsqu'on
l'évalue en $A$. Pour tout $A$-module $M$, nous pouvons choisir une suite exacte
$$
\bigoplus\nolimits_{j \in J} A \to
\bigoplus\nolimits_{i \in I} A \to
M \to 0
$$
En utilisant les morphismes construits ci-dessus, nous obtenons un diagramme commutatif
$$
\xymatrix{
(\bigoplus\nolimits_{j \in J} A) \otimes_A K \ar[r] \ar[d] &
(\bigoplus\nolimits_{i \in I} A) \otimes_A K \ar[r] \ar[d] &
M \otimes_A K \ar[r] \ar[d] &
0 \\
F(\bigoplus\nolimits_{j \in J} A) \ar[r] &
F(\bigoplus\nolimits_{i \in I} A) \ar[r] &
F(M) \ar[r] &  0
}
$$
La ligne inférieure est exacte puisque $F$ est exact à droite.
La ligne supérieure est exacte puisque le produit tensoriel par $K$ est exact à droite.
Comme $F$ commute aux sommes directes, les deux flèches verticales de gauche
sont bijectives. Nous pouvons donc conclure.
\end{proof}

\begin{example}
\label{example-functor-modules}
Soit $R$ un anneau. Soient $A$ et $B$ des $R$-algèbres. Soit $K$ un
$A \otimes_R B$-module. Nous pouvons alors considérer le foncteur
\begin{equation}
\label{equation-FM-modules}
F : \text{Mod}_A \longrightarrow \text{Mod}_B,\quad
M \longmapsto M \otimes_A K
\end{equation}
Ce foncteur est $R$-linéaire, exact à droite,
commute aux sommes directes quelconques, commute
à toutes les limites inductives et admet un adjoint à droite (lemme \ref{lemma-functor}).
\end{example}

\begin{lemma}
\label{lemma-functor-modules}
Soit $R$ un anneau. Soient $A$ et $B$ des $R$-algèbres. Il existe une
équivalence de catégories entre
\begin{enumerate}
\item la catégorie des foncteurs $R$-linéaires
$F : \text{Mod}_A \to \text{Mod}_B$ qui
sont exacts à droite et commutent aux sommes directes quelconques, et
\item la catégorie $\text{Mod}_{A \otimes_R B}$,
\end{enumerate}
donnée en associant à $K$ le foncteur $F$ de (\ref{equation-FM-modules}).
\end{lemma}

\begin{proof}
Soit $F$ un objet de la première catégorie. D'après le
lemme \ref{lemma-functor}, nous pouvons supposer que $F(M) = M \otimes_A K$
fonctoriellement en $M$, pour un certain $A \otimes_\mathbf{Z} B$-module $K$.
La $R$-linéarité de $F$ implique immédiatement que la structure de
$A \otimes_\mathbf{Z} B$-module sur $K$ provient
d'une (unique) structure de $A \otimes_R B$-module sur $K$.
Nous voyons donc que l'association de $K$ à $F$ comme dans (\ref{equation-FM-modules})
est essentiellement surjective.

\medskip\noindent
Pour démontrer que notre foncteur est pleinement fidèle, nous devons montrer que,
étant donnés des $A \otimes_R B$-modules $K$ et $K'$, toute transformation
$t : F \to F'$ entre les foncteurs correspondants provient
d'un unique $\varphi : K \to K'$. Puisque $K = F(A)$ et $K' = F'(A)$,
nous pouvons prendre pour $\varphi$ la valeur $t_A : F(A) \to F'(A)$
de $t$ en $A$. Cet homomorphisme est $A \otimes_R B$-linéaire d'après la
définition de la structure de $A \otimes B$-module sur $F(A)$
et $F'(A)$ donnée dans la démonstration du lemme \ref{lemma-functor}.
\end{proof}

\begin{remark}
\label{remark-composition}
Soit $R$ un anneau. Soient $A$, $B$, $C$ des $R$-algèbres.
Soient $F : \text{Mod}_A \to \text{Mod}_B$ et
$F' : \text{Mod}_B \to \text{Mod}_C$ des foncteurs
$R$-linéaires exacts à droite qui commutent aux sommes directes quelconques.
Si, par l'équivalence du lemme \ref{lemma-functor-modules}, l'objet
$K$ de $\text{Mod}_{A \otimes_R B}$ correspond à $F$ et l'objet
$K'$ de $\text{Mod}_{B \otimes_R C}$ correspond à $F'$, alors
$K \otimes_B K'$, vu comme un objet de
$\text{Mod}_{A \otimes_R C}$, correspond à $F' \circ F$.
\end{remark}

\begin{remark}
\label{remark-exact-flat}
Dans la situation du lemme \ref{lemma-functor-modules},
supposons que $F$ corresponde à $K$. Alors
$F$ est exact $\Leftrightarrow$ $K$ est plat sur $A$.
\end{remark}

\begin{remark}
\label{remark-finite}
Dans la situation du lemme \ref{lemma-functor-modules},
supposons que $F$ corresponde à $K$. Alors
$F$ envoie les $A$-modules finis sur des $B$-modules finis
$\Leftrightarrow$ $K$ est fini comme $B$-module.
\end{remark}

\begin{remark}
\label{remark-finite-presentation}
Dans la situation du lemme \ref{lemma-functor-modules},
supposons que $F$ corresponde à $K$. Alors
$F$ envoie les $A$-modules de présentation finie sur des $B$-modules de présentation finie
$\Leftrightarrow$ $K$ est de présentation finie comme $B$-module.
\end{remark}

\begin{lemma}
\label{lemma-functor-equivalence}
Soient $A$ et $B$ des anneaux. Si
$$
F : \text{Mod}_A \longrightarrow \text{Mod}_B
$$
est une équivalence de catégories, alors il existe un isomorphisme
d'anneaux $A \to B$ et un $B$-module inversible $L$ tels que
$F$ soit isomorphe au foncteur $M \mapsto (M \otimes_A B) \otimes_B L$.
\end{lemma}

\begin{proof}
Puisqu'une équivalence commute à toutes les limites inductives, nous voyons que le
lemme \ref{lemma-functor} s'applique. Soit $K$ le
$A \otimes_\mathbf{Z} B$-module tel que $F$ soit
isomorphe au foncteur $M \mapsto M \otimes_A K$.
Soit $K'$ le $B \otimes_\mathbf{Z} A$-module tel qu'un
quasi-inverse de $F$ soit
isomorphe au foncteur $N \mapsto N \otimes_B K'$.
D'après la remarque \ref{remark-composition} et le
lemme \ref{lemma-functor-modules}, nous avons un isomorphisme
$$
\psi : K \otimes_B K' \longrightarrow A
$$
de $A \otimes_\mathbf{Z} A$-modules.
De même, nous avons un isomorphisme
$$
\psi' : K' \otimes_A K \longrightarrow B
$$
de $B \otimes_\mathbf{Z} B$-modules. Choisissons un élément
$\xi = \sum_{i = 1, \ldots, n} x_i \otimes y_i \in K \otimes_B K'$
tel que $\psi(\xi) = 1$. Considérons les isomorphismes
$$
K \xrightarrow{\psi^{-1} \otimes \text{id}_K}
K \otimes_B K' \otimes_A K \xrightarrow{\text{id}_K \otimes \psi'} K
$$
Leur composé est un isomorphisme donné par
$$
k \longmapsto \sum x_i \psi'(y_i \otimes k)
$$
Nous en concluons que cet automorphisme se factorise comme
$$
K \to B^{\oplus n} \to K
$$
en tant qu'homomorphisme de $B$-modules. Il s'ensuit que $K$ est un
$B$-module projectif fini.

\medskip\noindent
Nous affirmons que $K$ est inversible comme $B$-module. Cela revient
à demander que le rang de $K$ comme $B$-module ait la valeur constante $1$ ;
voir Plus sur l'algèbre, lemme \ref{more-algebra-lemma-invertible}, et
Algèbre, lemme \ref{algebra-lemma-finite-projective}.
Sinon, il existe un idéal maximal $\mathfrak m \subset B$
tel que soit (a) $K \otimes_B B/\mathfrak m = 0$, soit
(b) il existe une surjection $K \to (B/\mathfrak m)^{\oplus 2}$ de
$B$-modules. Le cas (a) est absurde puisque $K' \otimes_A K \otimes_B N = N$
pour tout $B$-module $N$. Le cas (b) impliquerait l'existence d'une surjection
$$
A = K \otimes_B K' \longrightarrow (B/\mathfrak m \otimes_B K')^{\oplus 2}
$$
de $A$-modules (à droite). C'est impossible, car la cible est un $A$-module
qui nécessite au moins deux générateurs : $B/\mathfrak m \otimes_B K'$
est non nul, car il est l'image du module non nul $B/\mathfrak m$ par
le quasi-inverse de $F$.

\medskip\noindent
Puisque $K$ est inversible comme $B$-module, nous avons $\Hom_B(K, K) = B$.
Comme $K = F(A)$, l'action de $A$ sur $K$ définit un isomorphisme d'anneaux
$A \to B$. Le lemme en résulte.
\end{proof}

\begin{lemma}
\label{lemma-functor-equivalence-linear}
Soit $R$ un anneau. Soient $A$ et $B$ des $R$-algèbres. Si
$$
F : \text{Mod}_A \longrightarrow \text{Mod}_B
$$
est une équivalence $R$-linéaire de catégories, alors il existe
un isomorphisme $A \to B$ de $R$-algèbres et un $B$-module inversible $L$ tels que
$F$ soit isomorphe au foncteur $M \mapsto (M \otimes_A B) \otimes_B L$.
\end{lemma}

\begin{proof}
Le lemme \ref{lemma-functor-equivalence} nous fournit $A \to B$ et $L$.
Pour achever la démonstration, il reste à montrer que la $R$-linéarité
de $F$ force $A \to B$ à être un homomorphisme de $R$-algèbres. Nous omettons les détails.
\end{proof}

\begin{remark}
\label{remark-monoidal}
Soient $A$ et $B$ des anneaux. Munissons $\text{Mod}_A$ et $\text{Mod}_B$
de la structure monoïdale usuelle donnée par les produits tensoriels de modules.
Soit $F : \text{Mod}_A \to \text{Mod}_B$ un foncteur de
catégories monoïdales, voir
Catégories, définition \ref{categories-definition-functor-monoidal-categories}.
Voici quelques remarques :
\begin{enumerate}
\item Puisque $F(A)$ est une unité (d'après nos définitions), nous avons $F(A) = B$.
\item Nous obtenons une application multiplicative $\varphi : A \to B$
en envoyant $a \in A$ sur son action sur $F(A) = B$.
\item Prenons $A = B$ et $F(M) = M \otimes_A M$. Dans ce cas, $\varphi(a) = a^2$.
\item Si $F$ est additif, alors $\varphi$ est un homomorphisme d'anneaux.
\item Prenons $A = B = \mathbf{Z}$ et $F(M) = M/\text{torsion}$. Alors
$\varphi = \text{id}_\mathbf{Z}$, mais $F$ n'est pas le foncteur identité.
\item Si $F$ est exact à droite et commute aux sommes directes,
alors $F(M) = M \otimes_{A, \varphi} B$ d'après le lemme \ref{lemma-functor}.
\end{enumerate}
Autrement dit, les homomorphismes d'anneaux $A \to B$ sont en bijection avec les classes d'isomorphisme
de foncteurs de catégories monoïdales $\text{Mod}_A \to \text{Mod}_B$
qui commutent à toutes les limites inductives.
\end{remark}




\section{Prolongement des foncteurs sur des catégories de modules}
\label{section-functors-extend}

\noindent
Pour un anneau $A$, notons $\text{Mod}^{fp}_A$ la catégorie des
$A$-modules de présentation finie.

\begin{lemma}
\label{lemma-functor-fp-modules}
Soient $A$ et $B$ des anneaux. Soit
$F : \text{Mod}^{fp}_A \to \text{Mod}^{fp}_B$ un foncteur.
Alors $F$ se prolonge de manière unique en un foncteur
$F' : \text{Mod}_A \to \text{Mod}_B$
qui commute aux limites inductives filtrantes.
\end{lemma}

\begin{proof}
C'est un cas particulier du lemme \ref{lemma-functor-on-fp-modules}.
\end{proof}

\begin{remark}
\label{remark-monoidal-extension}
Conservons $A$, $B$, $F$ et $F'$ comme dans le lemme \ref{lemma-functor-fp-modules}.
Observons que le produit tensoriel de deux modules de présentation finie est
de présentation finie ; voir Algèbre, lemme \ref{algebra-lemma-tensor-finiteness}.
Nous pouvons donc munir $\text{Mod}^{fp}_A$, $\text{Mod}^{fp}_B$,
$\text{Mod}_A$ et $\text{Mod}_B$ de la structure monoïdale usuelle
donnée par les produits tensoriels de modules. Dans ce cas, si $F$ est
un foncteur de catégories monoïdales, alors $F'$ l'est aussi. Cela résulte immédiatement
du fait que les produits tensoriels de modules commutent aux limites inductives
filtrantes.
\end{remark}

\begin{lemma}
\label{lemma-functor-fp-modules-exact}
Conservons $A$, $B$, $F$ et $F'$ comme dans le lemme \ref{lemma-functor-fp-modules}.
\begin{enumerate}
\item Si $F$ est additif, alors $F'$ est additif et
commute aux sommes directes quelconques ;
\item si $F$ est exact à droite, alors $F'$ est exact à droite.
\end{enumerate}
\end{lemma}

\begin{proof}
Cela résulte des lemmes \ref{lemma-functor-on-fp-modules-additive} et
\ref{lemma-functor-on-fp-modules-right-exact}.
\end{proof}

\begin{remark}
\label{remark-monoidal-extension-exact}
En combinant les remarques \ref{remark-monoidal} et \ref{remark-monoidal-extension}
et le lemme \ref{lemma-functor-fp-modules-exact},
nous obtenons ce qui suit. Étant donnés des anneaux $A$ et $B$, l'ensemble des homomorphismes d'anneaux $A \to B$
est en bijection avec l'ensemble des classes d'isomorphisme
de foncteurs de catégories monoïdales $\text{Mod}^{fp}_A \to \text{Mod}^{fp}_B$
qui sont exacts à droite.
\end{remark}

\begin{lemma}
\label{lemma-functor-fp-modules-left-exact}
Conservons $A$, $B$, $F$ et $F'$ comme dans le lemme \ref{lemma-functor-fp-modules}.
Supposons que $A$ soit un anneau cohérent
(Algèbre, définition \ref{algebra-definition-coherent}).
Si $F$ est exact à gauche, alors $F'$ est exact à gauche.
\end{lemma}

\begin{proof}
C'est un cas particulier du lemme \ref{lemma-functor-on-fp-modules-left-exact}.
\end{proof}

\noindent
Pour un anneau $A$, notons $\text{Mod}^{fg}_A$ la catégorie des
$A$-modules de type fini (aussi appelés $A$-modules finis).

\begin{lemma}
\label{lemma-functor-finite-modules}
Soient $A$ et $B$ des anneaux noethériens. Soit
$F : \text{Mod}^{fg}_A \to \text{Mod}^{fg}_B$ un foncteur.
Alors $F$ se prolonge de manière unique en un foncteur $F' : \text{Mod}_A \to \text{Mod}_B$
qui commute aux limites inductives filtrantes. Si $F$ est additif, alors
$F'$ est additif et commute aux sommes directes quelconques.
Si $F$ est exact, exact à gauche ou exact à droite, alors $F'$ l'est aussi.
\end{lemma}

\begin{proof}
Voir les lemmes \ref{lemma-functor-fp-modules-exact} et
\ref{lemma-functor-fp-modules-left-exact}.
Utilisons aussi le fait que les $A$-modules finis sont des $A$-modules de présentation finie ;
voir Algèbre, lemme
\ref{algebra-lemma-Noetherian-finite-type-is-finite-presentation},
ainsi que le fait que les anneaux noethériens sont cohérents ; voir
Algèbre, lemme \ref{algebra-lemma-Noetherian-coherent}.
\end{proof}









\section{Foncteurs entre catégories de modules quasi-cohérents}
\label{section-functor-quasi-coherent}

\noindent
Dans cette section, nous étudions brièvement les foncteurs entre catégories de
modules quasi-cohérents.

\begin{example}
\label{example-functor-quasi-coherent}
Soit $R$ un anneau. Soient $X$ et $Y$ des
schémas sur $R$, avec $X$ quasi-compact et quasi-séparé.
Soit $\mathcal{K}$ un $\mathcal{O}_{X \times_R Y}$-module quasi-cohérent.
Nous pouvons alors considérer le foncteur
\begin{equation}
\label{equation-FM-QCoh}
F : \QCoh(\mathcal{O}_X) \longrightarrow \QCoh(\mathcal{O}_Y),\quad
\mathcal{F} \longmapsto
\text{pr}_{2, *}(\text{pr}_1^*\mathcal{F}
\otimes_{\mathcal{O}_{X \times_R Y}} \mathcal{K})
\end{equation}
Le morphisme $\text{pr}_2$ est quasi-compact et quasi-séparé
(Schémas, lemmes \ref{schemes-lemma-quasi-compact-preserved-base-change}
et \ref{schemes-lemma-separated-permanence}). Par conséquent, l'image directe par
ce morphisme préserve les modules quasi-cohérents ; voir
Schémas, lemme \ref{schemes-lemma-push-forward-quasi-coherent}.
De plus, notre foncteur est $R$-linéaire et commute aux sommes directes quelconques ;
voir Cohomologie des schémas, lemme \ref{coherent-lemma-colimit-cohomology}.
\end{example}

\noindent
Le lemme suivant est une généralisation naturelle du
lemme \ref{lemma-functor-modules}.

\begin{lemma}
\label{lemma-functor-quasi-coherent-from-affine}
Soit $R$ un anneau. Soient $X$ et $Y$ des schémas sur $R$, avec $X$ affine.
Il existe une équivalence de catégories entre
\begin{enumerate}
\item la catégorie des foncteurs $R$-linéaires
$F : \QCoh(\mathcal{O}_X) \to \QCoh(\mathcal{O}_Y)$
qui sont exacts à droite et commutent aux sommes directes quelconques, et
\item la catégorie $\QCoh(\mathcal{O}_{X \times_R Y})$,
\end{enumerate}
donnée en associant à $\mathcal{K}$ le foncteur $F$ de (\ref{equation-FM-QCoh}).
\end{lemma}

\begin{proof}
Soit $\mathcal{K}$ un objet de $\QCoh(\mathcal{O}_{X \times_R Y})$
et soit $F_\mathcal{K}$ le foncteur (\ref{equation-FM-QCoh}). D'après la discussion de
l'exemple \ref{example-functor-quasi-coherent}, nous savons déjà que
$F$ est $R$-linéaire et commute aux sommes directes quelconques.
Comme $\text{pr}_2 : X \times_R Y \to Y$ est affine
(Morphismes, lemme \ref{morphisms-lemma-base-change-affine}), le foncteur
$\text{pr}_{2, *}$ est exact ; voir Cohomologie des schémas, lemme
\ref{coherent-lemma-relative-affine-vanishing}.
Ainsi, $F$ est aussi exact à droite ; autrement dit, $F$ est comme en (1).

\medskip\noindent
Soit $F$ comme en (1). Écrivons $X = \Spec(A)$. Considérons le
$\mathcal{O}_Y$-module quasi-cohérent $\mathcal{G} = F(\mathcal{O}_X)$.
Le foncteur $F$ induit une application $R$-linéaire
$A \to \text{End}_{\mathcal{O}_Y}(\mathcal{G})$,
$a \mapsto F(a \cdot \text{id})$. Ainsi, $\mathcal{G}$ est un module sur le faisceau d'anneaux
$$
A \otimes_R \mathcal{O}_Y = \text{pr}_{2, *}\mathcal{O}_{X \times_R Y}
$$
D'après Morphismes, lemme \ref{morphisms-lemma-affine-equivalence-modules},
nous trouvons qu'il existe un unique module quasi-cohérent $\mathcal{K}$
sur $X \times_R Y$ tel que $F(\mathcal{O}_X) = \mathcal{G} =
\text{pr}_{2, *}\mathcal{K}$, de manière compatible avec l'action
de $A$ et de $\mathcal{O}_Y$. Notons $F_\mathcal{K}$ le foncteur
donné par (\ref{equation-FM-QCoh}). Il existe une équivalence
$\text{Mod}_A \to \QCoh(\mathcal{O}_X)$ envoyant $A$ sur $\mathcal{O}_X$ ; voir
Schémas, lemme \ref{schemes-lemma-equivalence-quasi-coherent}.
Nous obtenons donc un isomorphisme $F \cong F_\mathcal{K}$ d'après le
lemme \ref{lemma-functor-on-modules}, car nous avons un isomorphisme
$F(\mathcal{O}_X) \cong F_\mathcal{K}(\mathcal{O}_X)$ compatible avec
l'action de $A$ par construction.

\medskip\noindent
Cela montre que le foncteur qui associe à $\mathcal{K}$ le foncteur $F_\mathcal{K}$
est essentiellement surjectif. Nous omettons la vérification de la pleine fidélité.
\end{proof}

\begin{remark}
\label{remark-affine-morphism}
Nous utiliserons ci-dessous que, pour un morphisme affine
$h : T \to S$, nous avons $h_*\mathcal{G} \otimes_{\mathcal{O}_S} \mathcal{H} =
h_*(\mathcal{G} \otimes_{\mathcal{O}_T} h^*\mathcal{H})$ pour
$\mathcal{G} \in \QCoh(\mathcal{O}_T)$ et
$\mathcal{H} \in \QCoh(\mathcal{O}_S)$. Cela résulte
immédiatement de la traduction en algèbre.
\end{remark}

\begin{lemma}
\label{lemma-functor-quasi-coherent-from-affine-compose}
Dans le lemme \ref{lemma-functor-quasi-coherent-from-affine}, supposons que $F$
corresponde à $\mathcal{K}$ dans $\QCoh(\mathcal{O}_{X \times_R Y})$.
Nous avons :
\begin{enumerate}
\item Si $f : X' \to X$ est un morphisme affine, alors $F \circ f_*$
correspond à $(f \times \text{id}_Y)^*\mathcal{K}$.
\item Si $g : Y' \to Y$ est un morphisme plat, alors $g^* \circ F$ correspond à
$(\text{id}_X \times g)^*\mathcal{K}$.
\item Si $j : V \to Y$ est une immersion ouverte, alors $j^* \circ F$
correspond à $\mathcal{K}|_{X \times_R V}$.
\end{enumerate}
\end{lemma}

\begin{proof}
Démonstration de (1). Considérons le diagramme commutatif
$$
\xymatrix{
X' \times_R Y \ar[rrd]^{\text{pr}'_2} \ar[rd]_{f \times \text{id}_Y}
\ar[dd]_{\text{pr}'_1} \\
& X \times_R Y \ar[r]_{\text{pr}_2} \ar[d]_{\text{pr}_1} & Y \\
X' \ar[r]^f & X
}
$$
Soit $\mathcal{F}'$ un module quasi-cohérent sur $X'$. Nous avons
\begin{align*}
\text{pr}_{2, *}(\text{pr}_1^*f_*\mathcal{F}'
\otimes_{\mathcal{O}_{X \times_R Y}} \mathcal{K})
& =
\text{pr}_{2, *}((f \times \text{id}_Y)_*
(\text{pr}'_1)^*\mathcal{F}'
\otimes_{\mathcal{O}_{X \times_R Y}} \mathcal{K}) \\
& =
\text{pr}_{2, *}(f \times \text{id}_Y)_*
\left((\text{pr}'_1)^*\mathcal{F}'
\otimes_{\mathcal{O}_{X' \times_R Y}}
(f \times \text{id}_Y)^*\mathcal{K})\right)  \\
& =
\text{pr}'_{2, *}((\text{pr}'_1)^*\mathcal{F}'
\otimes_{\mathcal{O}_{X' \times_R Y}} (f \times \text{id}_Y)^*\mathcal{K})
\end{align*}
Ici, la première égalité est le changement de base affine pour le carré
de gauche du diagramme ; voir
Cohomologie des schémas, lemme \ref{coherent-lemma-affine-base-change}.
La deuxième égalité résulte de la remarque \ref{remark-affine-morphism}.
La troisième égalité est la fonctorialité des images directes de modules.
Cela démontre (1).

\medskip\noindent
Démonstration de (2). Considérons le diagramme commutatif
$$
\xymatrix{
X \times_R Y' \ar[rr]_-{\text{pr}'_2} \ar[rd]^{\text{id}_X \times g}
\ar[rdd]_{\text{pr}'_1} & & Y' \ar[d]^g \\
& X \times_R Y \ar[r]_-{\text{pr}_2} \ar[d]^{\text{pr}_1} & Y \\
& X
}
$$
Nous avons
\begin{align*}
g^*\text{pr}_{2, *}(\text{pr}_1^*\mathcal{F}
\otimes_{\mathcal{O}_{X \times_R Y}} \mathcal{K})
& =
\text{pr}'_{2, *}(
(\text{id}_X \times g)^*(
\text{pr}_1^*\mathcal{F} \otimes_{\mathcal{O}_{X \times_R Y}} \mathcal{K})) \\
& =
\text{pr}'_{2, *}((\text{pr}'_1)^*\mathcal{F}
\otimes_{\mathcal{O}_{X \times_R Y'}}
(\text{id}_X \times g)^*\mathcal{K})
\end{align*}
La première égalité résulte du changement de base plat pour le carré du diagramme ; voir
Cohomologie des schémas, lemme \ref{coherent-lemma-flat-base-change-cohomology}.
La deuxième égalité résulte de la fonctorialité de l'image inverse et du fait que
l'image inverse d'un produit tensoriel est le produit tensoriel des images inverses.

\medskip\noindent
La partie (3) est un cas particulier de (2).
\end{proof}

\begin{lemma}
\label{lemma-functor-quasi-coherent-from-affine-diagonal-pre}
Soit $R$ un anneau. Soient $X$ et $Y$ des schémas sur $R$. Supposons que $X$
soit quasi-compact à diagonale affine. Soit
$F : \QCoh(\mathcal{O}_X) \to \QCoh(\mathcal{O}_Y)$
un foncteur $R$-linéaire exact à droite qui commute
aux sommes directes quelconques. Nous pouvons alors construire
\begin{enumerate}
\item un module quasi-cohérent $\mathcal{K}$ sur $X \times_R Y$, et
\item une transformation naturelle $t : F \to F_\mathcal{K}$,
où $F_\mathcal{K}$ désigne le foncteur (\ref{equation-FM-QCoh}),
\end{enumerate}
telle que $t : F \circ f_* \to F_\mathcal{K} \circ f_*$ soit un isomorphisme
pour tout morphisme $f : X' \to X$ dont le schéma source est affine.
\end{lemma}

\begin{proof}
Considérons un morphisme $f' : X' \to X$ avec $X'$ affine. Comme la
diagonale de $X$ est affine, nous voyons que $f'$ est un morphisme affine
(Morphismes, lemme \ref{morphisms-lemma-affine-permanence}).
Ainsi, $f'_* : \QCoh(\mathcal{O}_{X'}) \to \QCoh(\mathcal{O}_X)$
est un foncteur exact $R$-linéaire
(Cohomologie des schémas, lemme \ref{coherent-lemma-relative-affine-vanishing})
qui commute aux sommes directes
(Cohomologie des schémas, lemme \ref{coherent-lemma-colimit-cohomology}).
Ainsi, $F \circ f'_*$ est un foncteur $R$-linéaire exact à droite qui commute
aux sommes directes quelconques. Par conséquent,
$F \circ f'_* = F_{\mathcal{K}'}$ pour un certain $\mathcal{K}'$
sur $X' \times_R Y$, d'après le lemme \ref{lemma-functor-quasi-coherent-from-affine}.
De plus, étant donné un morphisme $f'' : X'' \to X'$ avec $X''$ affine,
nous obtenons une identification canonique
$(f'' \times \text{id}_Y)^*\mathcal{K}' = \mathcal{K}''$
à l'aide des références déjà données, combinées au
lemme \ref{lemma-functor-quasi-coherent-from-affine-compose}.
Ces identifications satisfont une condition de cocycle pour
un autre morphisme $f''' : X''' \to X''$, condition dont nous laissons
au lecteur le soin d'écrire les détails.

\medskip\noindent
Choisissons un recouvrement ouvert affine $X = \bigcup_{i = 1, \ldots, n} U_i$.
Comme la diagonale de $X$ est affine, nous voyons que les intersections
$U_{i_0 \ldots i_p} = U_{i_0} \cap \ldots \cap U_{i_p}$ sont affines.
Comme ci-dessus, les morphismes d'inclusion
$j_{i_0 \ldots i_p} : U_{i_0 \ldots i_p} \to X$ sont affines.
Notons $\mathcal{K}_{i_0 \ldots i_p}$ le module quasi-cohérent
sur $U_{i_0 \ldots i_p} \times_R Y$ correspondant à
$F \circ j_{i_0 \ldots i_p *}$ comme ci-dessus.
D'après ce qui précède, nous obtenons des identifications
$$
\mathcal{K}_{i_0 \ldots i_p} =
\mathcal{K}_{i_0 \ldots \hat i_j \ldots i_p}|_{U_{i_0 \ldots i_p} \times_R Y}
$$
qui satisfont les compatibilités usuelles de recollement. Autrement dit, nous obtenons
un unique module quasi-cohérent $\mathcal{K}$ sur $X \times_R Y$
dont la restriction à $U_{i_0 \ldots i_p} \times_R Y$ est
$\mathcal{K}_{i_0 \ldots i_p}$, de manière compatible avec les identifications affichées.

\medskip\noindent
Construisons ensuite la transformation $t$. Étant donné un
$\mathcal{O}_X$-module quasi-cohérent $\mathcal{F}$, notons $\mathcal{F}_{i_0 \ldots i_p}$
la restriction de $\mathcal{F}$ à $U_{i_0 \ldots i_p}$ et notons
$(\text{pr}_1^*\mathcal{F} \otimes \mathcal{K})_{i_0 \ldots i_p}$
la restriction de $\text{pr}_1^*\mathcal{F} \otimes \mathcal{K}$ à
$U_{i_0 \ldots i_p} \times_R Y$.
Observons que
\begin{align*}
F(j_{i_0 \ldots i_p *}\mathcal{F}_{i_0 \ldots i_p})
& =
\text{pr}_{i_0 \ldots i_p, 2, *}(
\text{pr}_{i_0 \ldots i_p, 1}^*\mathcal{F}_{i_0 \ldots i_p}
\otimes \mathcal{K}_{i_0 \ldots i_p}) \\
& =
\text{pr}_{i_0 \ldots i_p, 2, *}
(\text{pr}_1^*\mathcal{F} \otimes \mathcal{K})_{i_0 \ldots i_p}
\end{align*}
où $\text{pr}_{i_0 \ldots i_p, 2} : U_{i_0 \ldots i_p} \times_R Y \to Y$
est la projection, et de même pour l'autre projection. De plus, ces
identifications sont compatibles avec les identifications affichées
au paragraphe précédent. Rappelons, d'après Cohomologie des schémas, lemme
\ref{coherent-lemma-separated-case-relative-cech},
que le complexe de {\v C}ech relatif
$$
\bigoplus 
\text{pr}_{i_0, 2, *}
(\text{pr}_1^*\mathcal{F} \otimes \mathcal{K})_{i_0}
\to
\bigoplus 
\text{pr}_{i_0i_1, 2, *}
(\text{pr}_1^*\mathcal{F} \otimes \mathcal{K})_{i_0i_1}
\to
\bigoplus 
\text{pr}_{i_0i_1i_2, 2, *}
(\text{pr}_1^*\mathcal{F} \otimes \mathcal{K})_{i_0i_1i_2}
\to \ldots
$$
calcule $R\text{pr}_{2, *}(\text{pr}_1^*\mathcal{F} \otimes \mathcal{K})$.
Ainsi, le faisceau de cohomologie en degré $0$ est $F_\mathcal{K}(\mathcal{F})$.
Nous obtenons donc le morphisme souhaité
$t : F(\mathcal{F}) \to F_\mathcal{K}(\mathcal{F})$
en considérant le diagramme commutatif suivant
$$
\xymatrix{
&
F(\mathcal{F}) \ar[r] \ar@{..>}[d] &
\bigoplus F(j_{i_0*}\mathcal{F}_{i_0}) \ar[r] \ar[d] &
\bigoplus F(j_{i_0i_1*}\mathcal{F}_{i_0i_1}) \ar[d] \\
0 \ar[r] &
F_\mathcal{K}(\mathcal{F}) \ar[r] &
\bigoplus 
\text{pr}_{i_0, 2, *}
(\text{pr}_1^*\mathcal{F} \otimes \mathcal{K})_{i_0}
\ar[r] &
\bigoplus 
\text{pr}_{i_0i_1, 2, *}
(\text{pr}_1^*\mathcal{F} \otimes \mathcal{K})_{i_0i_1}
}
$$
Nous obtenons la ligne supérieure en appliquant $F$ au complexe (exact)
$0 \to \mathcal{F} \to \bigoplus j_{i_0*}\mathcal{F}_{i_0} \to
\bigoplus j_{i_0i_1*}\mathcal{F}_{i_0i_1}$ (mais, comme
$F$ n'est pas exact, la ligne supérieure est seulement un complexe et n'est pas
nécessairement exacte).
Les flèches verticales pleines sont les identifications ci-dessus.
Cela définit bien la flèche pointillée comme souhaité.
La flèche est fonctorielle en $\mathcal{F}$ ; nous omettons les détails.

\medskip\noindent
Il nous reste à démontrer la dernière assertion. Soit $f : X' \to X$
comme dans l'énoncé du lemme, et soit $\mathcal{K}'$
le module quasi-cohérent sur $X' \times_R Y$ construit
au premier paragraphe de la démonstration. Si le morphisme
$f : X' \to X$ est à valeurs dans l'un des ouverts $U_i$, alors le
résultat découle du
lemme \ref{lemma-functor-quasi-coherent-from-affine-compose},
car, dans ce cas, nous savons
que $\mathcal{K}_i = \mathcal{K}|_{U_i \times_R Y}$
a pour image inverse $\mathcal{K}$. En général, nous obtenons un
recouvrement ouvert affine $X' = \bigcup U'_i$ avec $U'_i = f^{-1}(U_i)$,
ainsi que des isomorphismes
$\mathcal{K}'|_{U'_i} = f_i^*\mathcal{K}_i$, où
$f_i : U'_i \to U_i$ est le morphisme induit.
Ces morphismes satisfont les conditions de compatibilité nécessaires
pour se recoller en un isomorphisme $\mathcal{K}' = f^*\mathcal{K}$,
et nous concluons. Nous omettons certains détails.
\end{proof}

\begin{lemma}
\label{lemma-coh-noetherian-from-affine-flat}
Dans le lemme \ref{lemma-functor-quasi-coherent-from-affine}
ou dans le lemme \ref{lemma-functor-quasi-coherent-from-affine-diagonal-pre},
si $F$ est un foncteur exact, alors l'objet correspondant
$\mathcal{K}$ de $\QCoh(\mathcal{O}_{X \times_R Y})$ est plat sur $X$.
\end{lemma}

\begin{proof}
Nous pouvons supposer que $X$ est affine ; nous sommes donc dans le cas du
lemme \ref{lemma-functor-quasi-coherent-from-affine}.
D'après le lemme \ref{lemma-functor-quasi-coherent-from-affine-compose},
nous pouvons supposer que $Y$ est affine. Dans le cas affine, l'énoncé
se traduit par la remarque \ref{remark-exact-flat}.
\end{proof}

\begin{lemma}
\label{lemma-functor-quasi-coherent-from-affine-diagonal}
Soit $R$ un anneau. Soient $X$ et $Y$ des schémas sur $R$. Supposons que $X$ soit
quasi-compact à diagonale affine.
Il existe une équivalence de catégories entre
\begin{enumerate}
\item la catégorie des foncteurs exacts $R$-linéaires
$F : \QCoh(\mathcal{O}_X) \to \QCoh(\mathcal{O}_Y)$
qui commutent aux sommes directes quelconques, et
\item la sous-catégorie pleine de $\QCoh(\mathcal{O}_{X \times_R Y})$ constituée
des $\mathcal{K}$ tels que
\begin{enumerate}
\item $\mathcal{K}$ est plat sur $X$ ;
\item pour $\mathcal{F} \in \QCoh(\mathcal{O}_X)$, nous avons
$R^q\text{pr}_{2, *}(\text{pr}_1^*\mathcal{F}
\otimes_{\mathcal{O}_{X \times_R Y}} \mathcal{K}) = 0$ pour $q > 0$.
\end{enumerate}
\end{enumerate}
donnée en associant à $\mathcal{K}$ le foncteur $F$ de (\ref{equation-FM-QCoh}).
\end{lemma}

\begin{proof}
Soit $\mathcal{K}$ comme en (2). Le foncteur $F$ de
(\ref{equation-FM-QCoh}) commute aux sommes directes.
Puisque, d'après (1) (a), le module $\mathcal{K}$ est plat sur $X$,
nous voyons que, pour toute suite exacte courte
$0 \to \mathcal{F}_1 \to \mathcal{F}_2 \to \mathcal{F}_3 \to 0$,
nous obtenons une suite exacte courte
$$
0 \to
\text{pr}_1^*\mathcal{F}_1 \otimes_{\mathcal{O}_{X \times_R Y}} \mathcal{K} \to
\text{pr}_1^*\mathcal{F}_2 \otimes_{\mathcal{O}_{X \times_R Y}} \mathcal{K} \to
\text{pr}_1^*\mathcal{F}_3 \otimes_{\mathcal{O}_{X \times_R Y}} \mathcal{K} \to
0
$$
Puisque, d'après (2)(b), l'image directe supérieure $R^1\text{pr}_{2, *}$
du premier terme est nulle, nous concluons que
$0 \to F(\mathcal{F}_1) \to F(\mathcal{F}_2) \to F(\mathcal{F}_3) \to 0$
est exacte, et nous voyons que $F$ est comme en (1).

\medskip\noindent
Soit $F$ comme en (1). Soient $\mathcal{K}$ et $t : F \to F_\mathcal{K}$ comme
dans le lemme \ref{lemma-functor-quasi-coherent-from-affine-diagonal-pre}.
D'après le lemme \ref{lemma-coh-noetherian-from-affine-flat}, nous voyons
que $\mathcal{K}$ est plat sur $X$. Pour achever la démonstration, nous devons
montrer que $t$ est un isomorphisme et établir l'assertion sur les images
directes supérieures. Ces deux faits résultent de ce que le
complexe de {\v C}ech relatif
$$
\bigoplus 
\text{pr}_{i_0, 2, *}
(\text{pr}_1^*\mathcal{F} \otimes \mathcal{K})_{i_0}
\to
\bigoplus 
\text{pr}_{i_0i_1, 2, *}
(\text{pr}_1^*\mathcal{F} \otimes \mathcal{K})_{i_0i_1}
\to
\bigoplus 
\text{pr}_{i_0i_1i_2, 2, *}
(\text{pr}_1^*\mathcal{F} \otimes \mathcal{K})_{i_0i_1i_2}
\to \ldots
$$
calcule $R\text{pr}_{2, *}(\text{pr}_1^*\mathcal{F} \otimes \mathcal{K})$.
Voir la démonstration du
lemme \ref{lemma-functor-quasi-coherent-from-affine-diagonal-pre}
pour les notations et la raison de cette assertion. Dans la démonstration du
lemme \ref{lemma-functor-quasi-coherent-from-affine-diagonal-pre},
nous avons aussi constaté que ce complexe est égal à l'image par $F$ du complexe
$$
\bigoplus j_{i_0*}\mathcal{F}_{i_0} \to
\bigoplus j_{i_0i_1*}\mathcal{F}_{i_0i_1} \to
\bigoplus j_{i_0i_1i_2*}\mathcal{F}_{i_0i_1i_2} \to \ldots
$$
Ce complexe est exact sauf en degré zéro, où son faisceau de cohomologie
est égal à $\mathcal{F}$. Ainsi, puisque $F$ est un foncteur exact,
nous concluons que $F = F_\mathcal{K}$ et que (2)(b) est satisfaite.

\medskip\noindent
Nous omettons de démontrer que la construction qui associe à $F$
le module $\mathcal{K}$ est fonctorielle et est un quasi-inverse du
foncteur qui associe à $\mathcal{K}$ le foncteur $F_\mathcal{K}$
déterminé par (\ref{equation-FM-QCoh}).
\end{proof}

\begin{remark}
\label{remark-characterize-FM-QCoh}
Soit $R$ un anneau. Soient $X$ et $Y$ des schémas sur $R$. Supposons que $X$
soit quasi-compact à diagonale affine.
Le lemme \ref{lemma-functor-quasi-coherent-from-affine-diagonal} peut
se généraliser comme suit : les foncteurs
(\ref{equation-FM-QCoh}) associés aux modules quasi-cohérents sur
$X \times_R Y$ sont exactement ceux
$F : \QCoh(\mathcal{O}_X) \to \QCoh(\mathcal{O}_Y)$
qui possèdent les propriétés suivantes :
\begin{enumerate}
\item $F$ est $R$-linéaire et commute aux sommes directes quelconques ;
\item $F \circ j_*$ est exact à droite lorsque $j : U \to X$ est
l'inclusion d'un ouvert affine ;
\item $0 \to F(\mathcal{F}) \to F(\mathcal{G}) \to F(\mathcal{H})$
est exacte chaque fois que $0 \to \mathcal{F} \to \mathcal{G} \to \mathcal{H} \to 0$
est une suite exacte telle que, pour tout $x \in X$, la suite des germes
$0 \to \mathcal{F}_x \to \mathcal{G}_x \to \mathcal{H}_x \to 0$
soit une suite exacte courte scindée.
\end{enumerate}
En effet, ces hypothèses suffisent pour construire une transformation
$t : F \to F_\mathcal{K}$ comme dans le
lemme \ref{lemma-functor-quasi-coherent-from-affine-diagonal-pre}
et montrer qu'elle est un isomorphisme. De plus, les propriétés (1), (2) et (3)
sont bien satisfaites par les foncteurs (\ref{equation-FM-QCoh}).
Si nous en avons un jour besoin, nous l'énoncerons et le démontrerons ici avec soin.
\end{remark}

\begin{lemma}
\label{lemma-compose-FM-QCoh}
Soit $R$ un anneau. Soient $X$, $Y$, $Z$ des schémas sur $R$. Supposons que
$X$ et $Y$ soient quasi-compacts et à diagonale affine. Soient
$$
F : \QCoh(\mathcal{O}_X) \to \QCoh(\mathcal{O}_Y)
\quad\text{et}\quad
G : \QCoh(\mathcal{O}_Y) \to \QCoh(\mathcal{O}_Z)
$$
des foncteurs exacts $R$-linéaires qui commutent aux sommes directes quelconques.
Soient $\mathcal{K}$ dans $\QCoh(\mathcal{O}_{X \times_R Y})$
et $\mathcal{L}$ dans $\QCoh(\mathcal{O}_{Y \times_R Z})$
les « noyaux » correspondants ; voir le
lemme \ref{lemma-functor-quasi-coherent-from-affine-diagonal}.
Alors $G \circ F$ correspond à
$\text{pr}_{13, *}(\text{pr}_{12}^*\mathcal{K}
\otimes_{\mathcal{O}_{X \times_R Y \times_R Z}}
\text{pr}_{23}^*\mathcal{L})$ dans $\QCoh(\mathcal{O}_{X \times_R Z})$.
\end{lemma}

\begin{proof}
Puisque $G \circ F : \QCoh(\mathcal{O}_X) \to \QCoh(\mathcal{O}_Z)$
est $R$-linéaire, exact et commute aux sommes directes quelconques,
le lemme \ref{lemma-functor-quasi-coherent-from-affine-diagonal} montre
qu'il existe un objet $\mathcal{M}$ de
$\QCoh(\mathcal{O}_{X \times_R Z})$ correspondant à $G \circ F$.
D'autre part, notons
$\mathcal{E} = \text{pr}_{13, *}(\text{pr}_{12}^*\mathcal{K}
\otimes \text{pr}_{23}^*\mathcal{L})$. Ici et dans la suite de
la démonstration, nous omettons l'indice des produits tensoriels.
Soient $U \subset X$ et $W \subset Z$ des sous-schémas ouverts affines.
Pour démontrer le lemme, nous allons construire un isomorphisme
$$
\Gamma(U \times_R W, \mathcal{E})
\cong
\Gamma(U \times_R W, \mathcal{M})
$$
compatible avec les applications de restriction lorsque $U$ et $W$ varient.

\medskip\noindent
Tout d'abord, nous observons que
$$
\Gamma(U \times_R W, \mathcal{E}) =
\Gamma(U \times_R Y \times_R W,
\text{pr}_{12}^*\mathcal{K} \otimes \text{pr}_{23}^*\mathcal{L})
$$
par construction. Il nous faut donc montrer qu'il en va de même
pour $\mathcal{M}$.

\medskip\noindent
Écrivons $U = \Spec(A)$ et notons $j : U \to X$ le morphisme d'inclusion.
Rappelons, d'après la construction de $\mathcal{M}$ dans la démonstration du
lemme \ref{lemma-functor-quasi-coherent-from-affine}, que
$$
\Gamma(U \times_R W, \mathcal{M}) =
\Gamma(W, G(F(j_*\mathcal{O}_U)))
$$
où la structure de $A$-module du membre de droite est induite par
l'action de $A$ sur $\mathcal{O}_U$. La correspondance entre
$F$ et $\mathcal{K}$ nous dit que
$F(j_*\mathcal{O}_U) = b_*(a^*j_*\mathcal{O}_U \otimes \mathcal{K})$,
où $a : X \times_R Y \to X$ et $b : X \times_R Y \to Y$ sont
les morphismes de projection. Comme $j$ est un morphisme affine, nous avons
$a^*j_*\mathcal{O}_U = (j \times \text{id}_Y)_*\mathcal{O}_{U \times_R Y}$
d'après Cohomologie des schémas, lemme
\ref{coherent-lemma-affine-base-change}.
Ensuite, nous avons
$(j \times \text{id}_Y)_*\mathcal{O}_{U \times_R Y} \otimes \mathcal{K} =
(j \times \text{id}_Y)_*\mathcal{K}|_{U \times_R Y}$,
par exemple d'après la remarque \ref{remark-affine-morphism}.
En rassemblant ces égalités, nous obtenons
$$
F(j_*\mathcal{O}_U) =
(U \times_R Y \to Y)_*\mathcal{K}|_{U \times_R Y}
$$
muni de l'action évidente de $A$. (Cette formule est implicite dans la
démonstration du lemme \ref{lemma-functor-quasi-coherent-from-affine}.)
En appliquant le foncteur $G$, nous obtenons
$$
G(F(j_*\mathcal{O}_U)) =
t_*(s^*((U \times_R Y \to Y)_*\mathcal{K}|_{U \times_R Y})
\otimes \mathcal{L})
$$
où $s : Y \times_R Z \to Y$ et $t : Y \times_R Z \to Z$ sont
les morphismes de projection. En utilisant de nouveau le changement de base
pour un morphisme affine (Cohomologie des schémas, lemme
\ref{coherent-lemma-affine-base-change}), cette fois pour le carré
$$
\xymatrix{
U \times_R Y \times_R Z \ar[r] \ar[d] & U \times_R Y \ar[d] \\
Y \times_R Z \ar[r] & Y
}
$$
nous obtenons
$$
s^*((U \times_R Y \to Y)_*\mathcal{K}|_{U \times_R Y}) =
(U \times_R Y \times_R Z \to Y \times_R Z)_*
\text{pr}_{12}^*\mathcal{K}|_{U \times_R Y \times_R Z}
$$
En utilisant encore la remarque \ref{remark-affine-morphism}, nous trouvons
\begin{align*}
(U \times_R Y \times_R Z \to Y \times_R Z)_*
\text{pr}_{12}^*\mathcal{K}|_{U \times_R Y \times_R Z}
\otimes \mathcal{L} \\
=
(U \times_R Y \times_R Z \to Y \times_R Z)_*
\left(\text{pr}_{12}^*\mathcal{K} \otimes
\text{pr}_{23}^*\mathcal{L}\right)|_{U \times_R Y \times_R Z}
\end{align*}
En appliquant à cette égalité le foncteur
$\Gamma(W, t_*(-)) = \Gamma(Y \times_R W, -)$, nous obtenons
\begin{align*}
\Gamma(U \times_R W, \mathcal{M})
& =
\Gamma(W, G(F(j_*\mathcal{O}_U))) \\
& =
\Gamma(Y \times_R W, (U \times_R Y \times_R Z \to Y \times_R Z)_*
(\text{pr}_{12}^*\mathcal{K} \otimes
\text{pr}_{23}^*\mathcal{L})|_{U \times_R Y \times_R Z}) \\
& =
\Gamma(U \times_R Y \times_R W,
\text{pr}_{12}^*\mathcal{K} \otimes \text{pr}_{23}^*\mathcal{L})
\end{align*}
comme voulu. Nous omettons la vérification que ces isomorphismes sont
compatibles avec les applications de restriction.
\end{proof}

\begin{lemma}
\label{lemma-persistence-exactness}
Soient $R$, $X$, $Y$ et $\mathcal{K}$ comme au point (2) du
lemme \ref{lemma-functor-quasi-coherent-from-affine-diagonal}.
Alors, pour tout schéma $T$ sur $R$, nous avons
$$
R^q\text{pr}_{13, *}(\text{pr}_{12}^*\mathcal{F}
\otimes_{\mathcal{O}_{T \times_R X \times_R Y}}
\text{pr}_{23}^*\mathcal{K}) = 0
$$
pour tout module quasi-cohérent $\mathcal{F}$ sur $T \times_R X$ et tout
$q > 0$.
\end{lemma}

\begin{proof}
La question est locale sur $T$ ; nous pouvons donc supposer $T$ affine.
Dans ce cas, nous pouvons considérer le diagramme
$$
\xymatrix{
T \times_R X \ar[d] &
T \times_R X \times_R Y \ar[d] \ar[l] \ar[r] &
T \times_R Y \ar[d] \\
X &
X \times_R Y \ar[l] \ar[r] &
Y
}
$$
dont les flèches verticales sont affines. En particulier, le foncteur image
directe par $T \times_R Y \to Y$ est fidèle et exact (Cohomologie des schémas,
lemme \ref{coherent-lemma-relative-affine-vanishing} et
Morphismes, lemme \ref{morphisms-lemma-affine-equivalence-modules}).
En parcourant le diagramme et en utilisant l'annulation des images directes
supérieures par les morphismes affines (voir la référence ci-dessus), nous
voyons qu'il suffit de démontrer que
$$
R^q\text{pr}_{2, *}(
\text{pr}_{23, *}(\text{pr}_{12}^*\mathcal{F}
\otimes_{\mathcal{O}_{T \times_R X \times_R Y}}
\text{pr}_{23}^*\mathcal{K})) =
R^q\text{pr}_{2, *}(
\text{pr}_{23, *}(\text{pr}_{12}^*\mathcal{F})
\otimes_{\mathcal{O}_{X \times_R Y}}
\mathcal{K}))
$$
est nul, ce qui résulte de l'hypothèse sur $\mathcal{K}$.
L'égalité découle de la remarque \ref{remark-affine-morphism}.
\end{proof}

\begin{lemma}
\label{lemma-functor-quasi-coherent-from-separated}
Dans le lemme \ref{lemma-functor-quasi-coherent-from-affine-diagonal},
supposons que $F$ et $\mathcal{K}$ se correspondent. Si $X$ est séparé et
plat sur $R$, alors il existe une surjection
$\mathcal{O}_X \boxtimes F(\mathcal{O}_X) \to \mathcal{K}$.
\end{lemma}

\begin{proof}
Soit $\Delta : X \to X \times_R X$ le morphisme diagonal et
posons $\mathcal{O}_\Delta = \Delta_*\mathcal{O}_X$.
Comme $\Delta$ est une immersion fermée, nous avons une suite exacte courte
$$
0 \to \mathcal{I} \to 
\mathcal{O}_{X \times_R X} \to \mathcal{O}_\Delta \to 0
$$
Comme $\mathcal{K}$ est plat sur $X$, son image inverse
$\text{pr}_{23}^*\mathcal{K}$ sur $X \times_R X \times_R Y$
est plate sur $X \times_R X$. Nous obtenons une suite exacte courte
$$
0 \to 
\text{pr}_{12}^*\mathcal{I}
\otimes
\text{pr}_{23}^*\mathcal{K} \to
\text{pr}_{23}^*\mathcal{K} \to
\text{pr}_{12}^*\mathcal{O}_\Delta
\otimes
\text{pr}_{23}^*\mathcal{K} \to 0
$$
sur $X \times_R X \times_R Y$ ; voir
Modules, lemme \ref{modules-lemma-pullback-tensor-flat-module}.
Ainsi, d'après le lemme \ref{lemma-persistence-exactness},
nous obtenons une surjection
$$
\text{pr}_{13, *}(\text{pr}_{23}^*\mathcal{K})
\to
\text{pr}_{13, *}(
\text{pr}_{12}^*\mathcal{O}_\Delta
\otimes
\text{pr}_{23}^*\mathcal{K})
$$
D'après le changement de base plat (Cohomologie des schémas, lemme
\ref{coherent-lemma-flat-base-change-cohomology}), la source de cette flèche
est égale à $\text{pr}_2^*\text{pr}_{2, *}\mathcal{K}
= \mathcal{O}_X \boxtimes F(\mathcal{O}_X)$. D'autre part, le but est
égal à
$$
\text{pr}_{13, *}(
\text{pr}_{12}^*\mathcal{O}_\Delta
\otimes
\text{pr}_{23}^*\mathcal{K}) =
\text{pr}_{13, *} (\Delta \times \text{id}_Y)_* \mathcal{K} =
\mathcal{K}
$$
ce qui achève la démonstration. La première égalité résulte par exemple de
Cohomologie, lemme \ref{cohomology-lemma-projection-formula-closed-immersion}
et du fait que $\text{pr}_{12}^*\mathcal{O}_\Delta =
(\Delta \times \text{id}_Y)_*\mathcal{O}_{X \times_R Y}$.
\end{proof}









\section{Reconstruction de Gabriel--Rosenberg}
\label{section-gabriel}

\noindent
Le titre de cette section renvoie à des résultats tels que la
proposition \ref{proposition-gabriel-rosenberg}.
Outre l'article original de Gabriel \cite{Gabriel}, on pourra consulter
\cite{Brandenburg}, qui démontre le résultat pour les schémas quasi-séparés
et examine la littérature. Dans cette section, nous démontrerons seulement
la reconstruction de Gabriel--Rosenberg pour les schémas quasi-compacts et
quasi-séparés.

\begin{lemma}
\label{lemma-categorically-compact-QCoh}
Soit $X$ un schéma quasi-compact et quasi-séparé.
Soit $\mathcal{F}$ un $\mathcal{O}_X$-module quasi-cohérent.
Alors $\mathcal{F}$ est un objet compact au sens catégorique de
$\QCoh(\mathcal{O}_X)$ si et seulement si $\mathcal{F}$ est de
présentation finie.
\end{lemma}

\begin{proof}
Voir Catégories, définition \ref{categories-definition-compact-object},
pour notre notion d'objet compact au sens catégorique dans une catégorie.
Si $\mathcal{F}$ est de présentation finie, alors elle est compacte au sens
catégorique d'après Modules, lemme
\ref{modules-lemma-finite-presentation-quasi-compact-colimit}.
Réciproquement, tout module quasi-cohérent $\mathcal{F}$ peut s'écrire
comme une limite inductive filtrante $\mathcal{F} = \colim \mathcal{F}_i$
de $\mathcal{O}_X$-modules de présentation finie (donc quasi-cohérents) ;
voir Propriétés, lemme
\ref{properties-lemma-directed-colimit-finite-presentation}.
Si $\mathcal{F}$ est compact au sens catégorique, il existe un indice $i$
et un morphisme $\mathcal{F} \to \mathcal{F}_i$ qui est inverse à droite du
morphisme donné $\mathcal{F}_i \to \mathcal{F}$.
Nous concluons que $\mathcal{F}$ est facteur direct d'un module de
présentation finie, et est donc lui-même de présentation finie.
\end{proof}

\begin{lemma}
\label{lemma-supported-on-support-pre}
Soit $X$ un schéma affine. Soit $\mathcal{F}$ un
$\mathcal{O}_X$-module de présentation finie. Soit $\mathcal{E}$ un
$\mathcal{O}_X$-module quasi-cohérent non nul. Si
$\text{Supp}(\mathcal{E}) \subset \text{Supp}(\mathcal{F})$,
alors il existe un morphisme non nul $\mathcal{F} \to \mathcal{E}$.
\end{lemma}

\begin{proof}
Traduisons l'énoncé en algèbre. Soit $A$ un anneau. Soit $M$ un
$A$-module de présentation finie. Soit $N$ un $A$-module non nul. Supposons
$\text{Supp}(N) \subset \text{Supp}(M)$. Il s'agit de montrer que
$\Hom_A(M, N)$ est non nul.
Nous pouvons supposer $N = A/I$ cyclique (en remplaçant $N$ par n'importe
lequel de ses sous-modules cycliques non nuls). Choisissons une présentation
$$
A^{\oplus m} \xrightarrow{T} A^{\oplus n} \to M \to 0
$$
Rappelons que $\text{Supp}(M)$ est défini par $\text{Fit}_0(M)$, qui est
l'idéal engendré par les mineurs $n \times n$ de la matrice $T$.
Voir Plus sur l'algèbre, lemme
\ref{more-algebra-lemma-fitting-ideal-basics}.
L'hypothèse $\text{Supp}(N) \subset \text{Supp}(M)$ signifie alors que
les éléments de $\text{Fit}_0(M)$ sont nilpotents dans $A/I$.
Considérons la suite exacte
$$
0 \to \Hom_A(M, A/I) \to (A/I)^{\oplus n} \xrightarrow{T^t} (A/I)^{\oplus m}
$$
Nous devons montrer que $T^t$ ne peut être injective ; nous invitons le
lecteur à en trouver une démonstration à partir de la nilpotence des éléments
de $\text{Fit}_0(M)$ dans $A/I$. Voici la nôtre.
Comme $\text{Fit}_0(M)$ est de type fini, cette nilpotence implique que
l'annulateur $J \subset A/I$ de $\text{Fit}_0(M)$ dans $A/I$ est non nul.
Pour montrer que $T^t$ n'est pas injective, nous pouvons localiser en un idéal
premier. En choisissant convenablement cet idéal premier, nous pouvons
supposer $A$ local et $J$ toujours non nul. Alors $T^t$ a un noyau non nul
d'après Plus sur l'algèbre, lemme
\ref{more-algebra-lemma-exact-length-1}.
\end{proof}

\begin{lemma}
\label{lemma-supported-on-support}
Soit $X$ un schéma quasi-compact et quasi-séparé.
Soit $\mathcal{F}$ un $\mathcal{O}_X$-module de présentation finie.
Les deux sous-catégories suivantes de $\QCoh(\mathcal{O}_X)$ sont égales :
\begin{enumerate}
\item la sous-catégorie pleine $\mathcal{A} \subset \QCoh(\mathcal{O}_X)$
dont les objets sont les modules quasi-cohérents dont le support est contenu,
au sens ensembliste, dans $\text{Supp}(\mathcal{F})$ ;
\item la plus petite sous-catégorie de Serre
$\mathcal{B} \subset \QCoh(\mathcal{O}_X)$ qui contient
$\mathcal{F}$ et qui est stable par extensions et par sommes directes
quelconques.
\end{enumerate}
\end{lemma}

\begin{proof}
Observons que l'énoncé a un sens puisque les $\mathcal{O}_X$-modules de
présentation finie sont quasi-cohérents.
Comme $\mathcal{A}$ est une sous-catégorie de Serre stable par extensions et
par sommes directes et comme $\mathcal{F}$ est un objet de $\mathcal{A}$,
nous voyons que $\mathcal{B} \subset \mathcal{A}$. Il reste donc à montrer
que $\mathcal{A}$ est contenue dans $\mathcal{B}$.

\medskip\noindent
Soit $\mathcal{E}$ un objet de $\mathcal{A}$. Il existe un sous-module
maximal $\mathcal{E}' \subset \mathcal{E}$ appartenant à $\mathcal{B}$.
En effet, soit $\mathcal{E}_i \subset \mathcal{E}$, $i \in I$, la famille
des sous-objets qui sont des objets de $\mathcal{B}$. Alors
$\bigoplus \mathcal{E}_i$ appartient à $\mathcal{B}$, de même que
$$
\mathcal{E}' = \Im(\bigoplus \mathcal{E}_i \longrightarrow \mathcal{E})
$$
C'est manifestement le sous-module maximal recherché.

\medskip\noindent
Supposons maintenant qu'il existe un morphisme non nul
$\mathcal{G} \to \mathcal{E}/\mathcal{E}'$
avec $\mathcal{G}$ dans $\mathcal{B}$. Alors
$\mathcal{G}' = \mathcal{E} \times_{\mathcal{E}/\mathcal{E}'} \mathcal{G}$
appartient à $\mathcal{B}$ : $\mathcal{E}'$ en est le sous-objet et
$\mathcal{G}$ le quotient. L'image de $\mathcal{G}' \to \mathcal{E}$ serait alors
strictement plus grande que $\mathcal{E}'$, ce qui contredirait la maximalité
de $\mathcal{E}'$. Il suffit donc de démontrer l'assertion du paragraphe
suivant.

\medskip\noindent
Soit $\mathcal{E}$ un objet non nul de $\mathcal{A}$. Nous affirmons qu'il
existe un morphisme non nul $\mathcal{G} \to \mathcal{E}$ avec
$\mathcal{G}$ dans $\mathcal{B}$. Nous allons le démontrer par récurrence sur
le nombre minimal $n$ d'ouverts affines $U_i$ de $X$ tels que
$\text{Supp}(\mathcal{E}) \subset U_1 \cup \ldots \cup U_n$.
Posons $U = U_n$ et notons $j : U \to X$ le morphisme d'inclusion.
Notons $\mathcal{E}' = \Im(\mathcal{E} \to j_*\mathcal{E}|_U)$.
Le noyau $\mathcal{E}''$ de la surjection
$\mathcal{E} \to \mathcal{E}'$ a son support contenu dans
$U_1 \cup \ldots \cup U_{n - 1}$. Ainsi, si $\mathcal{E}''$ est non nul,
le résultat découle de l'hypothèse de récurrence. Autrement dit, nous pouvons
supposer $\mathcal{E} \subset j_*\mathcal{E}|_U$.
En particulier, $\mathcal{E}|_U$ est non nul.
D'après le lemme \ref{lemma-supported-on-support-pre}, il existe une
morphisme non nul $\mathcal{F}|_U \to \mathcal{E}|_U$.
Il correspond à un morphisme
$$
\varphi : \mathcal{F} \longrightarrow j_*(\mathcal{E}|_U)
$$
dont la restriction à $U$ est non nulle.
En posant $\mathcal{G} = \varphi^{-1}(\mathcal{E})$, nous concluons.
\end{proof}

\begin{lemma}
\label{lemma-quotient-supported-on-closed}
Soit $X$ un schéma quasi-compact et quasi-séparé.
Soit $Z \subset X$ une partie fermée telle que $U = X \setminus Z$
soit quasi-compact. Soit $\mathcal{A} \subset \QCoh(\mathcal{O}_X)$
la sous-catégorie pleine dont les objets sont les modules quasi-cohérents
à support dans $Z$. Alors le foncteur de restriction
$\QCoh(\mathcal{O}_X) \to \QCoh(\mathcal{O}_U)$ induit
une équivalence $\QCoh(\mathcal{O}_X)/\mathcal{A} \cong \QCoh(\mathcal{O}_U)$.
\end{lemma}

\begin{proof}
D'après la propriété universelle de la construction quotient
(Homologie, lemme \ref{homology-lemma-serre-subcategory-is-kernel}),
nous obtenons bien un foncteur induit
$\QCoh(\mathcal{O}_X)/\mathcal{A} \cong \QCoh(\mathcal{O}_U)$.
Notons $j : U \to X$ le morphisme d'inclusion. Comme $j$ est quasi-compact
et quasi-séparé, nous obtenons un foncteur
$j_* : \QCoh(\mathcal{O}_U) \to \QCoh(\mathcal{O}_X)$.
Le lecteur vérifiera qu'il définit un quasi-inverse ; nous omettons les détails.
\end{proof}

\begin{lemma}
\label{lemma-characterize-affine}
Soit $X$ un schéma quasi-compact et quasi-séparé.
Si $\QCoh(\mathcal{O}_X)$ est équivalente à la catégorie des modules
sur un anneau, alors $X$ est affine.
\end{lemma}

\begin{proof}
Soit $F : \text{Mod}_R \to \QCoh(\mathcal{O}_X)$ une équivalence.
Alors $\mathcal{F} = F(R)$ possède les propriétés suivantes :
\begin{enumerate}
\item c'est un $\mathcal{O}_X$-module de présentation finie
(lemme \ref{lemma-categorically-compact-QCoh}) ;
\item $\Hom_X(\mathcal{F}, -)$ est exact ;
\item $\Hom_X(\mathcal{F}, \mathcal{F})$ est un anneau commutatif ;
\item tout objet de $\QCoh(\mathcal{O}_X)$ est quotient d'une somme directe
de copies de $\mathcal{F}$.
\end{enumerate}
Soit $x \in X$ un point fermé. Considérons la surjection
$$
\mathcal{O}_X \to i_*\kappa(x)
$$
où le but est l'image directe de $\kappa(x)$ par le morphisme d'inclusion
$i : x \to X$. Nous avons
$$
\Hom_X(\mathcal{F}, i_*\kappa(x)) =
\Hom_{\mathcal{O}_{X, x}}(\mathcal{F}_x, \kappa(x))
$$
Tout d'abord, (4) implique que $\mathcal{F}_x$ est non nul.
D'après (2), tout homomorphisme $\mathcal{F}_x \to \kappa(x)$ se relève en
un homomorphisme $\mathcal{F}_x \to \mathcal{O}_{X, x}$ (puisqu'il se
relève même en un morphisme global $\mathcal{F} \to \mathcal{O}_X$).
Comme $\mathcal{F}_x$ est un $\mathcal{O}_{X, x}$-module fini, il s'ensuit
que $\mathcal{F}_x$ est un $\mathcal{O}_{X, x}$-module libre non nul de rang fini.
Comme $\mathcal{F}$ est de présentation finie, cela implique que
$\mathcal{F}$ est libre de rang fini strictement positif sur un voisinage ouvert de $x$
(Modules, lemme \ref{modules-lemma-finite-presentation-stalk-free}).
Puisque toute partie fermée de $X$ contient un point fermé
(Topologie, lemme \ref{topology-lemma-quasi-compact-closed-point}),
il s'ensuit que $\mathcal{F}$ est localement libre de rang fini strictement positif.
De même, l'homomorphisme
$$
\Hom_X(\mathcal{F}, \mathcal{F}) \to
\Hom_X(\mathcal{F}, i_*i^*\mathcal{F}) =
\Hom_{\kappa(x)}(\mathcal{F}_x/\mathfrak m_x \mathcal{F}_x,
\mathcal{F}_x/\mathfrak m_x \mathcal{F}_x)
$$
est surjective. D'après la propriété (3), nous concluons que le rang de
$\mathcal{F}_x$ est nécessairement $1$. Ainsi $\mathcal{F}$ est un
$\mathcal{O}_X$-module inversible. Nous en déduisons que le foncteur
$$
\mathcal{H}
\longmapsto
\Gamma(X, \mathcal{H}) = \Hom_X(\mathcal{O}_X, \mathcal{H}) =
\Hom_X(\mathcal{F}, \mathcal{H} \otimes_{\mathcal{O}_X} \mathcal{F})
$$
sur $\QCoh(\mathcal{O}_X)$ est lui aussi exact. Il s'ensuit que le premier
groupe $\Ext$
$$
\Ext^1_{\QCoh(\mathcal{O}_X)}(\mathcal{O}_X, \mathcal{H}) = 0
$$
calculé dans la catégorie abélienne $\QCoh(\mathcal{O}_X)$ s'annule pour
tout $\mathcal{H}$ dans $\QCoh(\mathcal{O}_X)$. Or, comme
$\QCoh(\mathcal{O}_X) \subset \textit{Mod}(\mathcal{O}_X)$
est stable par extensions (Schémas, section
\ref{schemes-section-quasi-coherent}), nous voyons que le groupe $\Ext^1$
entre modules quasi-cohérents calculé dans $\QCoh(\mathcal{O}_X)$ est le même
que celui calculé dans $\textit{Mod}(\mathcal{O}_X)$. Nous concluons donc que
$$
H^1(X, \mathcal{H}) =
\Ext^1_{\textit{Mod}(\mathcal{O}_X)}(\mathcal{O}_X, \mathcal{H}) = 0
$$
pour tout $\mathcal{H}$ dans $\QCoh(\mathcal{O}_X)$.
Il en résulte que $X$ est affine, par exemple d'après
Cohomologie des schémas, lemme
\ref{coherent-lemma-quasi-compact-h1-zero-covering}.
\end{proof}

\begin{proposition}
\label{proposition-gabriel-rosenberg}
\begin{reference}
Cas particulier de \cite[Theorem 1.2]{Brandenburg}
\end{reference}
Soient $X$ et $Y$ des schémas quasi-compacts et quasi-séparés.
Si $F : \QCoh(\mathcal{O}_X) \to \QCoh(\mathcal{O}_Y)$
est une équivalence, alors il existe un isomorphisme de schémas
$f : Y \to X$ et un $\mathcal{O}_Y$-module inversible
$\mathcal{L}$ tels que $F(\mathcal{F}) = f^*\mathcal{F} \otimes \mathcal{L}$.
\end{proposition}

\begin{proof}
Bien entendu, $F$ est additif, exact, commute à toutes les limites,
à toutes les limites inductives, aux sommes directes, etc.
Soit $U \subset X$ un sous-schéma ouvert affine.
Soit $\mathcal{I} \subset \mathcal{O}_X$ un faisceau quasi-cohérent
d'idéaux de type fini tel que $Z = V(\mathcal{I})$ soit le complémentaire
de $U$ dans $X$ ; voir Propriétés, lemme
\ref{properties-lemma-quasi-coherent-finite-type-ideals}.
Alors $\mathcal{O}_X/\mathcal{I}$ est un $\mathcal{O}_X$-module de
présentation finie. Ainsi $\mathcal{G} = F(\mathcal{O}_X/\mathcal{I})$
est un $\mathcal{O}_Y$-module de présentation finie d'après le
lemme \ref{lemma-categorically-compact-QCoh}.
Notons $T \subset Y$ le support de $\mathcal{G}$ et posons
$V = Y \setminus T$. Comme $\mathcal{G}$ est de présentation finie,
le schéma $V$ est un ouvert quasi-compact de $Y$.
D'après le lemme \ref{lemma-supported-on-support}, nous voyons que $F$ induit
une équivalence entre
\begin{enumerate}
\item la sous-catégorie pleine de $\QCoh(\mathcal{O}_X)$ formée des modules
à support dans $Z$, et
\item la sous-catégorie pleine de $\QCoh(\mathcal{O}_Y)$ formée des modules
à support dans $T$.
\end{enumerate}
D'après le lemme \ref{lemma-quotient-supported-on-closed}, nous obtenons un
diagramme commutatif
$$
\xymatrix{
\QCoh(\mathcal{O}_X) \ar[r]_F \ar[d] &
\QCoh(\mathcal{O}_Y) \ar[d] \\
\QCoh(\mathcal{O}_U) \ar[r]^{F_U} &
\QCoh(\mathcal{O}_V)
}
$$
où les flèches verticales sont les foncteurs de restriction et les flèches
horizontales des équivalences. D'après le lemme
\ref{lemma-characterize-affine}, nous concluons que $V$ est affine.
Le cas affine est traité par le lemme \ref{lemma-functor-equivalence}.
Nous trouvons donc un isomorphisme $f_U : V \to U$ et un
$\mathcal{O}_V$-module inversible $\mathcal{L}_U$ tels que
$F_U$ soit le foncteur
$\mathcal{F} \mapsto f_U^*\mathcal{F} \otimes \mathcal{L}_U$.

\medskip\noindent
Pour achever la démonstration, il suffit d'observer que les diagrammes
ci-dessus satisfont une compatibilité évidente à l'égard des inclusions de
sous-schémas ouverts affines de $X$. Les morphismes $f_U$ et les modules
inversibles $\mathcal{L}_U$ se recollent donc. Nous omettons les détails.
\end{proof}









\section{Foncteurs entre catégories de modules cohérents}
\label{section-functor-coherent}

\noindent
Le lemme suivant garantit que nous pouvons utiliser les résultats sur les
foncteurs entre catégories de modules quasi-cohérents lorsque nous disposons
d'un foncteur entre catégories de modules cohérents.

\begin{lemma}
\label{lemma-functor-coherent}
Soient $X$ et $Y$ des schémas noethériens. Soit
$F : \textit{Coh}(\mathcal{O}_X) \to \textit{Coh}(\mathcal{O}_Y)$
un foncteur. Alors $F$ se prolonge de manière unique en un foncteur
$\QCoh(\mathcal{O}_X) \to \QCoh(\mathcal{O}_Y)$
qui commute aux limites inductives filtrantes.
Si $F$ est additif, son prolongement commute aux sommes directes quelconques.
Si $F$ est exact, exact à gauche ou exact à droite, il en va de même de son
prolongement.
\end{lemma}

\begin{proof}
L'existence et l'unicité du prolongement sont un fait général ; voir
Catégories, lemme \ref{categories-lemma-extend-functor-by-colim}.
Pour vérifier que ce lemme s'applique, observons que les modules cohérents
sont de présentation finie (Modules, lemme
\ref{modules-lemma-coherent-finite-presentation}) et sont donc des objets
compacts au sens catégorique de $\textit{Mod}(\mathcal{O}_X)$ d'après
Modules, lemme \ref{modules-lemma-finite-presentation-quasi-compact-colimit}.
Enfin, tout module quasi-cohérent est limite inductive filtrante de modules
cohérents, par exemple d'après Propriétés, lemme
\ref{properties-lemma-quasi-coherent-colimit-finite-type}.

\medskip\noindent
Supposons $F$ additif. Si
$\mathcal{F} = \bigoplus_{j \in J} \mathcal{H}_j$
avec les $\mathcal{H}_j$ quasi-cohérents, alors
$\mathcal{F} = \colim_{J' \subset J\text{ fini}}
\bigoplus_{j \in J'} \mathcal{H}_j$.
En notant encore $F$ le prolongement de $F$, nous obtenons
\begin{align*}
F(\mathcal{F})
& =
\colim_{J' \subset J\text{ fini}}
F(\bigoplus\nolimits_{j \in J'} \mathcal{H}_j) \\
& =
\colim_{J' \subset J\text{ fini}}
\bigoplus\nolimits_{j \in J'} F(\mathcal{H}_j) \\
& =
\bigoplus\nolimits_{j \in J} F(\mathcal{H}_j)
\end{align*}
Ainsi $F$ commute aux sommes directes quelconques.

\medskip\noindent
Supposons que
$0 \to \mathcal{F} \to \mathcal{F}' \to \mathcal{F}'' \to 0$
soit une suite exacte courte de $\mathcal{O}_X$-modules quasi-cohérents.
Écrivons $\mathcal{F}' = \bigcup \mathcal{F}'_i$ comme réunion de ses
sous-modules cohérents ; voir Propriétés, lemme
\ref{properties-lemma-quasi-coherent-colimit-finite-type}.
Notons $\mathcal{F}''_i \subset \mathcal{F}''$ l'image de $\mathcal{F}'_i$
et posons $\mathcal{F}_i = \mathcal{F} \cap \mathcal{F}'_i =
\Ker(\mathcal{F}'_i \to \mathcal{F}''_i)$. Il est alors clair que
$\mathcal{F} = \bigcup \mathcal{F}_i$ et
$\mathcal{F}'' = \bigcup \mathcal{F}''_i$, et que nous avons des suites
exactes courtes
$$
0 \to \mathcal{F}_i \to \mathcal{F}_i' \to \mathcal{F}_i'' \to 0
$$
Comme le prolongement commute aux limites inductives filtrantes, nous avons
$F(\mathcal{F}) = \colim_{i \in I} F(\mathcal{F}_i)$,
$F(\mathcal{F}') = \colim_{i \in I} F(\mathcal{F}'_i)$ et
$F(\mathcal{F}'') = \colim_{i \in I} F(\mathcal{F}''_i)$.
Puisque les limites inductives filtrantes sont exactes
(Modules, lemme \ref{modules-lemma-limits-colimits}), nous concluons que les
propriétés d'exactitude de $F$ se transmettent à son prolongement.
\end{proof}

\begin{lemma}
\label{lemma-equivalence-coherent}
Soient $X$ et $Y$ des schémas noethériens. Soit
$F : \textit{Coh}(\mathcal{O}_X) \to \textit{Coh}(\mathcal{O}_Y)$
une équivalence de catégories. Alors il existe un isomorphisme $f : Y \to X$
et un $\mathcal{O}_Y$-module inversible $\mathcal{L}$ tels que
$F(\mathcal{F}) = f^*\mathcal{F} \otimes \mathcal{L}$.
\end{lemma}

\begin{proof}
D'après le lemme \ref{lemma-functor-coherent}, nous obtenons un unique foncteur
$F' : \QCoh(\mathcal{O}_X) \to \QCoh(\mathcal{O}_Y)$ qui prolonge $F$.
Il en va de même pour un quasi-inverse de $F$ et, par unicité, nous concluons
que $F'$ est une équivalence. D'après la
proposition \ref{proposition-gabriel-rosenberg}, il existe un isomorphisme
$f : Y \to X$ et un $\mathcal{O}_Y$-module inversible $\mathcal{L}$
tels que $F'(\mathcal{F}) = f^*\mathcal{F} \otimes \mathcal{L}$.
Alors $f$ et $\mathcal{L}$ conviennent également pour $F$.
\end{proof}

\begin{remark}
\label{remark-equivalence-coherent-linear}
Dans le lemme \ref{lemma-equivalence-coherent}, si $X$ et $Y$ sont définis
sur un même anneau de base $R$ et si $F$ est $R$-linéaire, alors
l'isomorphisme $f$ est un morphisme de schémas sur $R$.
\end{remark}

\begin{lemma}
\label{lemma-characterize-finite}
Soit $f : V \to X$ un morphisme quasi-fini séparé de schémas noethériens.
S'il existe un $\mathcal{O}_V$-module cohérent $\mathcal{K}$ dont le support
est $V$, tel que $f_*\mathcal{K}$ soit cohérent et que
$R^qf_*\mathcal{K} = 0$, alors $f$ est fini.
\end{lemma}

\begin{proof}
D'après le théorème principal de Zariski, il existe une immersion ouverte
$j : V \to Y$ au-dessus de $X$ telle que $\pi : Y \to X$ soit fini ; voir
Plus sur les morphismes, lemme
\ref{more-morphisms-lemma-quasi-finite-separated-pass-through-finite}.
Comme $\pi$ est affine, le foncteur $\pi_*$ est exact et fidèle sur la
catégorie des $\mathcal{O}_X$-modules cohérents.
Nous voyons donc que $j_*\mathcal{K}$ est cohérent et que
$R^qj_*\mathcal{K}$ est nul pour $q > 0$.
Autrement dit, nous sommes ramenés au cas traité dans le paragraphe suivant.

\medskip\noindent
Supposons que $f$ soit une immersion ouverte. Nous pouvons remplacer $X$ par
l'adhérence schématique de $V$. Supposons $X \setminus V$ non vide afin
d'aboutir à une contradiction. Choisissons un point générique
$\xi \in X \setminus V$ d'une composante irréductible de $X \setminus V$.
En examinant la situation après le changement de base par
$\Spec(\mathcal{O}_{X, \xi}) \to X$, puis en utilisant le changement de base
plat et Cohomologie locale, lemme
\ref{local-cohomology-lemma-finiteness-pushforwards-and-H1-local},
nous sommes ramenés au problème algébrique du paragraphe suivant.

\medskip\noindent
Soit $(A, \mathfrak m)$ un anneau local noethérien. Soit $M$ un $A$-module
fini dont le support est $\Spec(A)$. Alors
$H^i_\mathfrak m(M) \not = 0$ pour un certain $i$. Cela résulte de
Complexes dualisants, lemme \ref{dualizing-lemma-depth}, et du fait que
$M$ est non nul, donc de profondeur finie.
\end{proof}

\noindent
Le lemme suivant se généralise au cas où $k$ est un anneau noethérien et où
$X$ est plat sur $k$ (toutes les autres hypothèses restant inchangées).

\begin{lemma}
\label{lemma-functor-coherent-over-field}
Soit $k$ un corps. Soient $X$ et $Y$ des schémas de type fini sur $k$,
avec $X$ séparé. Il existe une équivalence de catégories entre
\begin{enumerate}
\item la catégorie des foncteurs exacts $k$-linéaires
$F : \textit{Coh}(\mathcal{O}_X) \to \textit{Coh}(\mathcal{O}_Y)$, et
\item la catégorie des $\mathcal{O}_{X \times Y}$-modules cohérents
$\mathcal{K}$ qui sont plats sur $X$ et dont le support est fini sur $Y$,
\end{enumerate}
obtenue en associant à $\mathcal{K}$ la restriction du foncteur
(\ref{equation-FM-QCoh}) à $\textit{Coh}(\mathcal{O}_X)$.
\end{lemma}

\begin{proof}
Soit $\mathcal{K}$ comme en (2). D'après le
lemme \ref{lemma-functor-quasi-coherent-from-affine-diagonal}, le foncteur
$F$ donné par (\ref{equation-FM-QCoh}) est exact et $k$-linéaire.
De plus, $F$ envoie $\textit{Coh}(\mathcal{O}_X)$ dans
$\textit{Coh}(\mathcal{O}_Y)$, par exemple d'après
Cohomologie des schémas, lemme
\ref{coherent-lemma-support-proper-over-base-pushforward}.

\medskip\noindent
Construisons le quasi-inverse de cette construction. Soit $F$ comme en (1).
D'après le lemme \ref{lemma-functor-coherent}, nous pouvons prolonger $F$ en
un foncteur exact $k$-linéaire sur les catégories de modules quasi-cohérents,
qui commute aux sommes directes quelconques.
D'après le lemme \ref{lemma-functor-quasi-coherent-from-affine-diagonal}, ce
prolongement correspond à un unique module quasi-cohérent $\mathcal{K}$, plat
sur $X$, tel que
$R^q\text{pr}_{2, *}(\text{pr}_1^*\mathcal{F}
\otimes_{\mathcal{O}_{X \times Y}} \mathcal{K}) = 0$ pour $q > 0$
et pour tout $\mathcal{O}_X$-module quasi-cohérent $\mathcal{F}$.
Comme $F(\mathcal{O}_X)$ est un $\mathcal{O}_Y$-module cohérent, il résulte
du lemme \ref{lemma-functor-quasi-coherent-from-separated} que
$\mathcal{K}$ est cohérent.

\medskip\noindent
Pour tout point fermé $x \in X$, notons $\mathcal{O}_x$ le faisceau gratte-ciel
en $x$ de valeur le corps résiduel de $x$. Nous avons
$$
F(\mathcal{O}_x) =
\text{pr}_{2, *}(\text{pr}_1^*\mathcal{O}_x \otimes \mathcal{K}) =
(x \times Y \to Y)_*(\mathcal{K}|_{x \times Y})
$$
Comme $x \times Y \to Y$ est fini, l'image directe par ce morphisme est
fidèle. Ainsi, si $y \in Y$ appartient à l'image du support de
$\mathcal{K}|_{x \times Y}$, alors $y$ appartient au support de
$F(\mathcal{O}_x)$.

\medskip\noindent
Soit $Z \subset X \times Y$ le support schématique $Z$ de $\mathcal{K}$ ;
voir Morphismes, définition
\ref{morphisms-definition-scheme-theoretic-support}.
Nous montrons d'abord que $Z \to Y$ est quasi-fini en montrant que ses fibres
au-dessus des points fermés sont finies. En effet, si la fibre de $Z \to Y$
au-dessus d'un point fermé $y \in Y$ est de dimension $> 0$, nous pouvons
trouver une infinité de points fermés deux à deux distincts
$x_1, x_2, \ldots$ dans l'image de $Z_y \to X$. Comme nous avons une
surjection
$\mathcal{O}_X \to \bigoplus_{i = 1, \ldots, n} \mathcal{O}_{x_i}$,
nous obtenons une surjection
$$
F(\mathcal{O}_X) \to \bigoplus\nolimits_{i = 1, \ldots, n} F(\mathcal{O}_{x_i})
$$
D'après ce qui précède, le point $y$ appartient au support de chacun des
modules cohérents $F(\mathcal{O}_{x_i})$. Comme $F(\mathcal{O}_X)$ est un
module cohérent, cela conduit à une contradiction, car le germe de
$F(\mathcal{O}_X)$ en $y$ sera engendré par $< n$ éléments si
$n$ est assez grand.
Ainsi $Z \to Y$ est quasi-fini.
Comme $\text{pr}_{2, *}\mathcal{K}$ est cohérent et que
$R^q\text{pr}_{2, *}\mathcal{K} = 0$ pour $q > 0$, nous concluons que
$Z \to Y$ est fini d'après le lemme \ref{lemma-characterize-finite}.
\end{proof}

\begin{lemma}
\label{lemma-pushforward-invertible-pre}
Soit $f : X \to Y$ un morphisme séparé de type fini de schémas. Soit
$\mathcal{F}$ un module quasi-cohérent de type fini sur $X$, dont le support
est fini sur $Y$ et tel que $\mathcal{L} = f_*\mathcal{F}$ soit un
$\mathcal{O}_X$-module inversible. Alors il existe une section
$s : Y \to X$ telle que $\mathcal{F} \cong s_*\mathcal{L}$.
\end{lemma}

\begin{proof}
Localement dans le cas affine, l'énoncé se traduit par le problème algébrique
suivant. Soit $A \to B$ un homomorphisme d'anneaux et soit $N$ un $B$-module
qui est inversible comme $A$-module. Alors l'annulateur $J$ de $N$ dans $B$
est tel que $A \to B/J$ soit un isomorphisme. Nous omettons les détails.
\end{proof}

\begin{lemma}
\label{lemma-pushforward-invertible}
Soit $f : X \to Y$ un morphisme séparé de type fini de schémas, muni d'une
section $s : Y \to X$. Soit $\mathcal{F}$ un module quasi-cohérent de type
fini sur $X$, à support ensembliste dans $s(Y)$, et tel que
$\mathcal{L} = f_*\mathcal{F}$ soit un $\mathcal{O}_X$-module inversible.
Si $Y$ est réduit, alors $\mathcal{F} \cong s_*\mathcal{L}$.
\end{lemma}

\begin{proof}
D'après le lemme \ref{lemma-pushforward-invertible-pre}, il existe une section
$s' : Y  \to X$ telle que $\mathcal{F} = s'_*\mathcal{L}$. Comme $s'(Y)$
et $s(Y)$ ont la même partie fermée sous-jacente et sont tous deux des
sous-schémas fermés réduits de $X$, ils sont égaux.
Ainsi $s = s'$ et le lemme est démontré.
\end{proof}

\begin{lemma}
\label{lemma-equivalence-coherent-over-field}
\begin{reference}
Version faible du résultat de \cite{Gabriel} selon lequel la catégorie des
modules quasi-cohérents détermine la classe d'isomorphisme d'un schéma.
\end{reference}
Soit $k$ un corps. Soient $X$ et $Y$ des schémas de type fini sur $k$,
avec $X$ séparé et $Y$ réduit. S'il existe une équivalence $k$-linéaire
$F : \textit{Coh}(\mathcal{O}_X) \to \textit{Coh}(\mathcal{O}_Y)$
de catégories, alors il existe un isomorphisme $f : Y \to X$ sur $k$ et un
$\mathcal{O}_Y$-module inversible $\mathcal{L}$ tels que
$F(\mathcal{F}) = f^*\mathcal{F} \otimes \mathcal{L}$.
\end{lemma}

\begin{proof}[Démonstration par reconstruction de Gabriel--Rosenberg]
Ce lemme est une forme faible des résultats exposés dans le
lemme \ref{lemma-equivalence-coherent} et la
remarque \ref{remark-equivalence-coherent-linear}.
\end{proof}

\begin{proof}[Démonstration indépendante de la reconstruction de Gabriel--Rosenberg]
D'après le lemme \ref{lemma-functor-coherent-over-field}, nous obtenons un
$\mathcal{O}_{X \times Y}$-module cohérent $\mathcal{K}$, plat sur $X$ et
dont le support est fini sur $Y$, tel que $F$ soit donné par la restriction
du foncteur (\ref{equation-FM-QCoh}) à $\textit{Coh}(\mathcal{O}_X)$.
Si nous pouvons montrer que $F(\mathcal{O}_X)$ est un
$\mathcal{O}_Y$-module inversible, alors le
lemme \ref{lemma-pushforward-invertible-pre} montre que
$\mathcal{K} = s_*\mathcal{L}$ pour une section
$s : Y \to X \times Y$ de $\text{pr}_2$ et un
$\mathcal{O}_Y$-module inversible $\mathcal{L}$.
Cela montrera que $F$ est de la forme indiquée, avec
$f = \text{pr}_1 \circ s$. Nous omettons certains détails.

\medskip\noindent
Il reste à montrer que $F(\mathcal{O}_X)$ est inversible. Nous n'en donnons
qu'une esquisse et omettons certains détails.
Pour un point fermé $x \in X$, nous notons $\mathcal{O}_x$ dans
$\textit{Coh}(\mathcal{O}_X)$ le faisceau gratte-ciel en $x$ de valeur
$\kappa(x)$. Observons d'abord que les seuls objets simples de la catégorie
$\textit{Coh}(\mathcal{O}_X)$ sont ces faisceaux gratte-ciel
$\mathcal{O}_x$. Il en va de même pour $Y$. Ainsi, pour tout point fermé
$y \in Y$, il existe un point fermé $x \in X$ tel que
$\mathcal{O}_y \cong F(\mathcal{O}_x)$. De plus, en considérant les
endomorphismes, nous trouvons $\kappa(x) \cong \kappa(y)$ comme extensions
finies de $k$. Alors
$$
\Hom_Y(F(\mathcal{O}_X), \mathcal{O}_y) \cong
\Hom_Y(F(\mathcal{O}_X), F(\mathcal{O}_x)) \cong
\Hom_X(\mathcal{O}_X, \mathcal{O}_x) \cong \kappa(x) \cong \kappa(y)
$$
Il en résulte que le germe du $\mathcal{O}_Y$-module cohérent
$F(\mathcal{O}_X)$ en $y \in Y$ peut être engendré par $1$ générateur (et
pas moins), pour tout point fermé $y \in Y$. Il s'ensuit immédiatement que
$F(\mathcal{O}_X)$ est localement engendré par $1$ élément (et pas moins) et,
comme $Y$ est réduit, cela montre bien qu'il s'agit d'un module inversible.
\end{proof}









\begin{multicols}{2}[\section{Autres chapitres}]
\noindent
Préliminaires
\begin{enumerate}
\item \hyperref[introduction-section-phantom]{Introduction}
\item \hyperref[conventions-section-phantom]{Conventions}
\item \hyperref[sets-section-phantom]{Théorie des ensembles}
\item \hyperref[categories-section-phantom]{Catégories}
\item \hyperref[topology-section-phantom]{Topologie}
\item \hyperref[sheaves-section-phantom]{Faisceaux sur les espaces}
\item \hyperref[sites-section-phantom]{Sites et faisceaux}
\item \hyperref[stacks-section-phantom]{Champs}
\item \hyperref[fields-section-phantom]{Corps}
\item \hyperref[algebra-section-phantom]{Algèbre commutative}
\item \hyperref[brauer-section-phantom]{Groupes de Brauer}
\item \hyperref[homology-section-phantom]{Algèbre homologique}
\item \hyperref[derived-section-phantom]{Catégories dérivées}
\item \hyperref[simplicial-section-phantom]{Méthodes simpliciales}
\item \hyperref[more-algebra-section-phantom]{Compléments d'algèbre}
\item \hyperref[smoothing-section-phantom]{Lissage des morphismes d'anneaux}
\item \hyperref[modules-section-phantom]{Faisceaux de modules}
\item \hyperref[sites-modules-section-phantom]{Modules sur les sites}
\item \hyperref[injectives-section-phantom]{Objets injectifs}
\item \hyperref[cohomology-section-phantom]{Cohomologie des faisceaux}
\item \hyperref[sites-cohomology-section-phantom]{Cohomologie sur les sites}
\item \hyperref[dga-section-phantom]{Algèbre différentielle graduée}
\item \hyperref[dpa-section-phantom]{Algèbre à puissances divisées}
\item \hyperref[sdga-section-phantom]{Faisceaux différentiels gradués}
\item \hyperref[hypercovering-section-phantom]{Hyperrecouvrements}
\end{enumerate}
Schémas
\begin{enumerate}
\setcounter{enumi}{25}
\item \hyperref[schemes-section-phantom]{Schémas}
\item \hyperref[constructions-section-phantom]{Constructions de schémas}
\item \hyperref[properties-section-phantom]{Propriétés des schémas}
\item \hyperref[morphisms-section-phantom]{Morphismes de schémas}
\item \hyperref[coherent-section-phantom]{Cohomologie des schémas}
\item \hyperref[divisors-section-phantom]{Diviseurs}
\item \hyperref[limits-section-phantom]{Limites de schémas}
\item \hyperref[varieties-section-phantom]{Variétés}
\item \hyperref[topologies-section-phantom]{Topologies sur les schémas}
\item \hyperref[descent-section-phantom]{Descente}
\item \hyperref[perfect-section-phantom]{Catégories dérivées de schémas}
\item \hyperref[more-morphisms-section-phantom]{Compléments sur les morphismes}
\item \hyperref[flat-section-phantom]{Compléments sur la platitude}
\item \hyperref[groupoids-section-phantom]{Schémas en groupoïdes}
\item \hyperref[more-groupoids-section-phantom]{Compléments sur les schémas en groupoïdes}
\item \hyperref[etale-section-phantom]{Morphismes étales de schémas}
\end{enumerate}
Thèmes de théorie des schémas
\begin{enumerate}
\setcounter{enumi}{41}
\item \hyperref[chow-section-phantom]{Homologie de Chow}
\item \hyperref[intersection-section-phantom]{Théorie de l'intersection}
\item \hyperref[pic-section-phantom]{Schémas de Picard des courbes}
\item \hyperref[weil-section-phantom]{Théories cohomologiques de Weil}
\item \hyperref[adequate-section-phantom]{Modules adéquats}
\item \hyperref[dualizing-section-phantom]{Complexes dualisants}
\item \hyperref[duality-section-phantom]{Dualité pour les schémas}
\item \hyperref[discriminant-section-phantom]{Discriminants et différentes}
\item \hyperref[derham-section-phantom]{Cohomologie de de Rham}
\item \hyperref[local-cohomology-section-phantom]{Cohomologie locale}
\item \hyperref[algebraization-section-phantom]{Géométrie algébrique et formelle}
\item \hyperref[curves-section-phantom]{Courbes algébriques}
\item \hyperref[resolve-section-phantom]{Résolution des surfaces}
\item \hyperref[models-section-phantom]{Réduction semi-stable}
\item \hyperref[functors-section-phantom]{Foncteurs et morphismes}
\item \hyperref[equiv-section-phantom]{Catégories dérivées de variétés}
\item \hyperref[pione-section-phantom]{Groupes fondamentaux de schémas}
\item \hyperref[etale-cohomology-section-phantom]{Cohomologie étale}
\item \hyperref[crystalline-section-phantom]{Cohomologie cristalline}
\item \hyperref[proetale-section-phantom]{Cohomologie pro-étale}
\item \hyperref[relative-cycles-section-phantom]{Cycles relatifs}
\item \hyperref[more-etale-section-phantom]{Compléments de cohomologie étale}
\item \hyperref[trace-section-phantom]{La formule des traces}
\end{enumerate}
Espaces algébriques
\begin{enumerate}
\setcounter{enumi}{64}
\item \hyperref[spaces-section-phantom]{Espaces algébriques}
\item \hyperref[spaces-properties-section-phantom]{Propriétés des espaces algébriques}
\item \hyperref[spaces-morphisms-section-phantom]{Morphismes d'espaces algébriques}
\item \hyperref[decent-spaces-section-phantom]{Espaces algébriques décents}
\item \hyperref[spaces-cohomology-section-phantom]{Cohomologie des espaces algébriques}
\item \hyperref[spaces-limits-section-phantom]{Limites d'espaces algébriques}
\item \hyperref[spaces-divisors-section-phantom]{Diviseurs sur les espaces algébriques}
\item \hyperref[spaces-over-fields-section-phantom]{Espaces algébriques sur des corps}
\item \hyperref[spaces-topologies-section-phantom]{Topologies sur les espaces algébriques}
\item \hyperref[spaces-descent-section-phantom]{Descente et espaces algébriques}
\item \hyperref[spaces-perfect-section-phantom]{Catégories dérivées d'espaces}
\item \hyperref[spaces-more-morphisms-section-phantom]{Compléments sur les morphismes d'espaces}
\item \hyperref[spaces-flat-section-phantom]{Platitude sur les espaces algébriques}
\item \hyperref[spaces-groupoids-section-phantom]{Groupoïdes dans les espaces algébriques}
\item \hyperref[spaces-more-groupoids-section-phantom]{Compléments sur les groupoïdes dans les espaces}
\item \hyperref[bootstrap-section-phantom]{Amorçage}
\item \hyperref[spaces-pushouts-section-phantom]{Sommes amalgamées d'espaces algébriques}
\end{enumerate}
Thèmes de géométrie
\begin{enumerate}
\setcounter{enumi}{81}
\item \hyperref[spaces-chow-section-phantom]{Groupes de Chow des espaces}
\item \hyperref[groupoids-quotients-section-phantom]{Quotients de groupoïdes}
\item \hyperref[spaces-more-cohomology-section-phantom]{Compléments sur la cohomologie des espaces}
\item \hyperref[spaces-simplicial-section-phantom]{Espaces simpliciaux}
\item \hyperref[spaces-duality-section-phantom]{Dualité pour les espaces}
\item \hyperref[formal-spaces-section-phantom]{Espaces algébriques formels}
\item \hyperref[restricted-section-phantom]{Algébrisation des espaces formels}
\item \hyperref[spaces-resolve-section-phantom]{Résolution des surfaces revisitée}
\end{enumerate}
Théorie des déformations
\begin{enumerate}
\setcounter{enumi}{89}
\item \hyperref[formal-defos-section-phantom]{Théorie formelle des déformations}
\item \hyperref[defos-section-phantom]{Théorie des déformations}
\item \hyperref[cotangent-section-phantom]{Le complexe cotangent}
\item \hyperref[examples-defos-section-phantom]{Problèmes de déformation}
\end{enumerate}
Champs algébriques
\begin{enumerate}
\setcounter{enumi}{93}
\item \hyperref[algebraic-section-phantom]{Champs algébriques}
\item \hyperref[examples-stacks-section-phantom]{Exemples de champs}
\item \hyperref[stacks-sheaves-section-phantom]{Faisceaux sur les champs algébriques}
\item \hyperref[criteria-section-phantom]{Critères de représentabilité}
\item \hyperref[artin-section-phantom]{Axiomes d'Artin}
\item \hyperref[quot-section-phantom]{Espaces de Quot et de Hilbert}
\item \hyperref[stacks-properties-section-phantom]{Propriétés des champs algébriques}
\item \hyperref[stacks-morphisms-section-phantom]{Morphismes de champs algébriques}
\item \hyperref[stacks-limits-section-phantom]{Limites de champs algébriques}
\item \hyperref[stacks-cohomology-section-phantom]{Cohomologie des champs algébriques}
\item \hyperref[stacks-perfect-section-phantom]{Catégories dérivées de champs}
\item \hyperref[stacks-introduction-section-phantom]{Introduction aux champs algébriques}
\item \hyperref[stacks-more-morphisms-section-phantom]{Compléments sur les morphismes de champs}
\item \hyperref[stacks-geometry-section-phantom]{La géométrie des champs}
\end{enumerate}
Thèmes de théorie des espaces de modules
\begin{enumerate}
\setcounter{enumi}{107}
\item \hyperref[moduli-section-phantom]{Champs de modules}
\item \hyperref[moduli-curves-section-phantom]{Espaces de modules de courbes}
\end{enumerate}
Divers
\begin{enumerate}
\setcounter{enumi}{109}
\item \hyperref[examples-section-phantom]{Exemples}
\item \hyperref[exercises-section-phantom]{Exercices}
\item \hyperref[guide-section-phantom]{Guide bibliographique}
\item \hyperref[desirables-section-phantom]{Desiderata}
\item \hyperref[coding-section-phantom]{Style de programmation}
\item \hyperref[obsolete-section-phantom]{Obsolète}
\item \hyperref[fdl-section-phantom]{Licence de documentation libre GNU}
\item \hyperref[index-section-phantom]{Index généré automatiquement}
\end{enumerate}
\end{multicols}


\bibliography{my}
\bibliographystyle{amsalpha}

\end{document}
